% संपूर्ण मराठी ग्रंथातील प्रकरण 24: संगणनक्षमतेचा सिद्धांत
% हा संपादनयोग्य स्रोतखंड आहे. संकलनासाठी संपूर्ण स्रोतसंग्रहातील
% अक्षररूपे, पूर्वभाग, चिन्हव्याख्या आणि आकृती-स्रोत आवश्यक आहेत.
% मूळ: Open Logic Project; मजकूर आणि रूपांतर CC BY 4.0.
% यंत्रानुवाद, दुरुस्ती आणि तपासणी: OpenAI Codex —
% GPT-5.6 Sol आणि GPT-6 Sol, दोन्ही Ultra effort.
% दुरुस्ती-पुनरावलोकन: GPT-6.1 Sol, Ultra effort.
\chapter{संगणनक्षमतेचा सिद्धांत}\label{cmp:thy::chap}
\begin{editorial}
  या प्रकरणातील मजकुराचे पुनरावलोकन करून त्याचा विस्तार करणे आवश्यक आहे.
  विशेषतः, येथे अजून सराव दिलेले नाहीत.
\end{editorial}








\section{परिचय}\label{cmp:thy:int:sec}

\emph{संगणनक्षमता सिद्धांत} म्हणून ओळखली जाणारी तर्कशास्त्राची शाखा
फलने व संच यांची संगणनक्षमता किंवा सापेक्ष संगणनक्षमता यांच्याशी संबंधित
प्रश्नांचा अभ्यास करते. या विषयाला पूर्वी \emph{पुनरावर्तन सिद्धांत}
म्हणत असत, यावरून क्लिनींचा प्रभाव दिसतो; आज ही दोन्ही नावे
सर्वसाधारणपणे वापरली जातात.

एखाद्या संगणनाच्या प्रतिमानात फलन~$f\colon \Nat \pto \Nat$ चे संगणन
करता येत असेल, तर त्याला \emph{आंशिक संगणनक्षम} म्हणू. $f$ सर्वत्र
परिभाषित असेल, तर $f$ ला फक्त \emph{संगणनक्षम} म्हणू. ज्याचे
जातिबोधक फलन~$\Char{R}$ संगणनक्षम आहे, असा संबंध~$R$ देखील
संगणनक्षम म्हणतात. $f$ व~$g$ ही आंशिक फलने असतील, तर
$f(x) \fdefined$ याचा अर्थ $f$ हे~$x$ येथे परिभाषित आहे, म्हणजे
$x$ हे~$f$ च्या परिभाषाक्षेत्रात आहे; $f(x) \fundefined$ याचा
अर्थ उलट, म्हणजे $f$ हे~$x$ येथे परिभाषित नाही. $f(x) \simeq g(x)$
याचा अर्थ $f(x)$ आणि $g(x)$ ही दोन्ही अपरिभाषित आहेत, किंवा दोन्ही
परिभाषित असून समान आहेत.

संगणनाच्या एखाद्या विशिष्ट प्रतिमानाचा उल्लेख न करताही संगणनक्षमता
सिद्धांताचा अभ्यास करता येतो. त्यासाठी एक सार्वत्रिक आंशिक संगणनक्षम
फलन~$\fn{Un}(k, x)$ आहे, हे दाखवतात. त्याच्या आधारे आंशिक संगणनक्षम
फलनांचे प्रगणन करता येते. $k$-वे एकस्थानी आंशिक संगणनक्षम फलन दर्शवण्यासाठी
आपण~$\cfind{k}$ हे चिन्हांकन वापरू; त्याची व्याख्या
$\cfind{k}(x) \simeq \fn{Un}(k, x)$ अशी आहे. (क्लिनी यांनी यासाठी
$\{ k \}$ हे चिन्हांकन वापरले होते; पण अलीकडे ते तितकेसे वापरले
जात नाही.) यापेक्षा थोड्या अधिक सामान्य स्वरूपात, कोणत्याही
स्थानसंख्येच्या आंशिक संगणनक्षम फलनांचे एकसमान प्रगणन करता येते.
$k$-वे $n$-स्थानी आंशिक पुनरावर्ती फलन दर्शवण्यासाठी आपण
$\cfind{k}[n]$ वापरू.

$f(\vec x, y)$ हे सर्वत्र परिभाषित किंवा आंशिक फलन असेल, तर
$\umin{y}{f (\vec x, y)}$ हे~$\vec x$ चे असे फलन आहे जे
$f(\vec x, y) = 0$ होईल असा लघुतम~$y$ परत करते, मात्र
$f(\vec x, 0)$, \dots, $f(\vec x, y-1)$ ही सर्व मूल्ये परिभाषित
असली पाहिजेत. असा~$y$ नसल्यास $\umin{y}{f (\vec x, y)}$
अपरिभाषित असते. $R(\vec x, y)$ हा संबंध असेल, तर
$\umin{y}{R(\vec x, y)}$ ची व्याख्या $R(\vec x, y)$ सत्य
असेल असा लघुतम~$y$ अशी केली जाते; दुसऱ्या शब्दांत,
$1 \tsub \Char{R}(\vec x, y) = 0$ होईल असा लघुतम~$y$.

एखादे फलन संगणनक्षम आहे हे दाखवण्यासाठी पुढील दोन मार्ग आहेत:
\begin{enumerate}
\item काटेकोरपणे: ट्यूरिंग यंत्र किंवा आंशिक पुनरावर्ती फलन स्पष्टपणे
  वर्णन करणे आणि ते इच्छित फलनाचे संगणन करते हे दाखवणे;
\item अनौपचारिकपणे: त्याचे संगणन करणाऱ्या अल्गोरिदमचे वर्णन करणे आणि
  चर्च यांच्या प्रबंधाचा आधार घेणे.
\end{enumerate}
या दोन मार्गांमध्ये नेमकी सीमारेषा नाही. अल्गोरिदमचे तपशीलवार
वर्णन पुरेशी माहिती देणारे असावे, जेणेकरून तत्त्वतः योग्य ट्यूरिंग
यंत्र किंवा आंशिक पुनरावर्ती व्याख्यांची मालिका कशी तयार करायची हे
तुलनेने स्पष्ट होईल. पूर्णपणे काटेकोर व्याख्या फारशा माहितीपूर्ण
असण्याची शक्यता नाही. म्हणून आपण या दोन मार्गांमध्ये योग्य मध्यममार्ग
शोधू; थोडक्यात, अचूक व्याख्या मिळवता येतील याची खात्री देणारे
सहज समजणारे पण काटेकोर पुरावे शोधू.











\section{संगणनांचे सांकेतिकरण}\label{cmp:thy:cod:sec}

संगणनाच्या प्रत्येक प्रतिमानात पुढील गोष्टी करता येतात:
\begin{enumerate}
\item संगणनक्षम फलनांच्या \emph{व्याख्यांचे} पद्धतशीर वर्णन करता येते.
  उदाहरणार्थ, ट्यूरिंग यंत्रांचे विनिर्देश, पुनरावर्ती व्याख्या किंवा
  कार्यक्रमण भाषेतील कार्यक्रम अशा व्याख्या देतात असे मानता येते.
\item एखाद्या व्याख्येने दिलेल्या फलनाचे ठरावीक आदानासाठी झालेले
  संगणन पूर्णपणे नोंदवता येते. उदाहरणार्थ, संगणनाच्या प्रत्येक पायरीवरील
  विन्यासांची क्रमिका (यंत्राची अवस्था आणि पट्टीवरील मजकूर) ट्यूरिंग
  यंत्राचे संगणन वर्णन करते.
\item संगणनाची एखादी कथित नोंद खरोखरच, दिलेल्या व्याख्येच्या
  संगणनक्षम फलनाचे दिलेल्या आदानासाठी झालेले संगणन कसे घडले याची नोंद
  आहे का, हे तपासता येते (त्या आदानावर फलन परिभाषित आहे, म्हणजे
  संगणन थांबते).
\item संगणनाच्या अशा पूर्ण नोंदीच्या वर्णनातून दिलेल्या आदानासाठी
  फलनाचे मूल्य काढता येते. उदाहरणार्थ, थांबणाऱ्या ट्यूरिंग यंत्राच्या
  संगणनाच्या अगदी शेवटच्या पायरीत पट्टीवर असलेला मजकूर हेच मूल्य असते.
\end{enumerate}

सांकेतिकरण वापरून संगणनक्षम फलनाच्या प्रत्येक वर्णनाला संख्यात्मक
\emph{निर्देशांक} देता येतो, आणि त्या निर्देशांकावरून सूचना संगणनक्षम
रीतीने पुन्हा मिळवता येतात. त्याचप्रमाणे, संगणनाची पूर्ण नोंदही एका
संख्येने संकेतबद्ध करता येते. मग ``निर्देशांक~$e$ असलेल्या फलनाच्या
आदान~$x$ वरील संगणनाच्या नोंदीचा संकेतांक~$s$ आहे'' हा परिणामी
अंकगणितीय संबंध आणि ``संकेतांक~$s$ असलेल्या संगणनक्रमिकेचे निर्गत''
हे फलन संगणनक्षम असतात; किंबहुना ती आदिम पुनरावर्ती असतात.

हे मूलभूत तथ्य फार सामर्थ्यवान आहे. त्याच्या आधारे संगणनक्षमतेविषयी
अनेक महत्त्वाचे आणि आश्चर्यकारक परिणाम सिद्ध करता येतात, तेही निवडलेल्या
संगणन-प्रतिमानावर अवलंबून न राहता.











\section{प्रमाण रूप प्रमेय}\label{cmp:thy:nfm:sec}

संगणनक्षम फलनांच्या व्याख्यांचे वर्णन करता येते आणि एखाद्या आदानावर
त्या फलनाच्या संगणनाची पूर्ण नोंद असल्याचे सांगणारे कथित वर्णन खरे आहे
का हे तपासता येते, असे समजा. मग संगणनाच्या प्रतिमानावर अवलंबून न राहता,
कोणत्याही संगणनक्षम फलन~$f$ चे कोणत्याही आदान~$x$ वरील मूल्य पुढीलप्रमाणे
ठरवता येईल, असे वाटते:
\begin{enumerate}
  \item संगणनांच्या नोंदींच्या सर्व शक्य वर्णनांमध्ये शोध घेणे.
  \item एखादी नोंद~$f(x)$ च्या संगणनाची नोंद आहे का हे तपासणे.
  \item योग्य नोंद सापडल्यास तिच्यातून~$f(x)$ चे मूल्य काढणे.
\end{enumerate}
हे खरोखरच शक्य आहे, हाच क्लिनी यांच्या प्रमाण रूप प्रमेयाचा आशय आहे.

\begin{thm}[क्लिनीचे प्रमाण रूप प्रमेय]
\label{cmp:thy:nfm:thm:normal-form}
असा आदिम पुनरावर्ती संबंध~$T(e, x, s)$ आणि असे आदिम पुनरावर्ती
फलन~$U(s)$ असते की पुढील गुणधर्म पूर्ण होतो: $f$ हे कोणतेही
आंशिक संगणनक्षम फलन असेल, तर एखाद्या~$e$ साठी
\[f(x) \simeq U(\umin{s}{T(e, x, s)})
\]
हे सर्व~$x$ साठी असते.
\end{thm}

\begin{proof}[सिद्धतेचे रूपरेषात्मक वर्णन]
संगणनाच्या कोणत्याही प्रतिमानासाठी संगणनक्षम फलन~$f$ चे वर्णन
काटेकोरपणे परिभाषित करून ते वर्णन नैसर्गिक संख्या~$e$ वापरून
संकेतबद्ध करता येते. त्याचप्रमाणे, निर्देशांक~$e$ असलेल्या फलनाचे
आदान~$x$ वरील संगणन कसे घडले याची नोंद करणाऱ्या ``संगणनक्रमिके''ची
कल्पना काटेकोरपणे परिभाषित करता येते. अशा संगणनक्रमिकेला देखील
संख्या~$s$ ने संकेतबद्ध करता येते. हे अशा प्रकारे करता येते की
\begin{enumerate}
  \item संख्या~$s$ ही निर्देशांक~$e$ असलेल्या फलनाच्या आदान~$x$
  वरील संगणनक्रमिकेचा संकेतांक असेल तेव्हा आणि केवळ तेव्हाच खरा असणारा संबंध
  $T(e, x, s)$, आणि
  \item $s$ ने संकेतबद्ध केलेल्या संगणनक्रमिकेला त्या संगणनाचा अंतिम
  परिणाम देणारे फलन $U(s)$
\end{enumerate}
ही दोन्ही संगणनक्षम असतात. किंबहुना, संबंध~$T$ आणि फलन~$U$
आदिम पुनरावर्ती असतात.
\end{proof}

\begin{explain}
प्रमाण रूप प्रमेयाची काटेकोर सिद्धता द्यायची असेल, तर संगणनाचे एक
प्रतिमान ठरवून संगणनक्षम फलनांच्या वर्णनांचे आणि संगणनक्रमिकांचे
सांकेतिकरण तपशीलवार करावे लागेल; तसेच संबंध~$T$ आणि फलन~$U$
आदिम पुनरावर्ती आहेत हे पडताळावे लागेल. बहुतेक उपयोगांसाठी $T$
आणि~$U$ संगणनक्षम आहेत व $U$ सर्वत्र परिभाषित आहे, एवढे पुरेसे असते.

प्रमाण रूप प्रमेयाची सिद्धता एका संक्षिप्त सूत्रवाक्यात लक्षात ठेवणे
कदाचित उत्तम: $\umin{s}{T(e, x, s)}$ हे निर्देशांक~$e$ असलेल्या
फलनाच्या आदान~$x$ वरील संगणनक्रमिकेचा शोध घेते; अशी क्रमिका
सापडल्यास $U$ त्या संगणनाचे निर्गत परत करते.
\end{explain}

संगणनाचे प्रतिमान आंशिक पुनरावर्ती फलने हे असेल (क्लिनी यांनी
मूळतः हेच प्रतिमान वापरले होते), तर या प्रमेयावरून कोणतेही फलन
परिभाषित करण्यासाठी अपरिबद्ध शोधाचा, म्हणजे एकाच
$\umin{y}{f(\vec x, y)}$ प्रवर्तकाचा, एकदाच वापर पुरेसा असतो.
या अर्थाने प्रत्येक आंशिक पुनरावर्ती फलनाला प्रमाण रूप असते.

$T$ आणि $U$ वापरून $\cfind{0}$, $\cfind{1}$, $\cfind{2}$, \dots
हे प्रगणन परिभाषित करता येते. यापुढे $T$ आणि~$U$ यांची एक योग्य
निवड निश्चित केली आहे असे समजू, आणि
\[\cfind{e}(x) \simeq U(\umin{s}{T(e,x,s)})
\]
हे समीकरण $\cfind{e}$ ची \emph{व्याख्या} मानू.

आणखी एक उपयुक्त तथ्य पुढीलप्रमाणे आहे:

\begin{thm}
प्रत्येक आंशिक संगणनक्षम फलनाला अनंत निर्देशांक असतात.
\end{thm}

हेही सहज स्पष्ट होते. एखाद्या संगणनक्षम फलनाचे वर्णन दिले असता,
त्याच फलनाचे (समान आदान-निर्गत जोड्यांचे) संगणन करणारे वेगळे वर्णन
तयार करता येते. उदाहरणार्थ, संगणनावर काहीही परिणाम न करणारी क्रिया
आधी करता येते ($0 = 0$ आहे का ते तपासणे किंवा $5$ पर्यंत मोजणे इत्यादी).
बदललेल्या वर्णनाचा निर्देशांक मूळ निर्देशांकापेक्षा नेहमी वेगळा असेल.
दोन्ही त्याच फलनाचे निर्देशांक आहेत; फक्त संगणनाची पद्धत किंचित वेगळी आहे.











\section{$s$-$m$-$n$ प्रमेय}\label{cmp:thy:smn:sec}

\begin{explain}
पुढील प्रमेयाला ``$s$-$m$-$n$ प्रमेय'' म्हणतात; त्याचे कारण थोड्याच
वेळात स्पष्ट होईल. हे प्रमेय नेमके काय सांगते हे समजून घेणे हाच कठीण
भाग आहे; विधान समजल्यावर ते बरेचसे स्पष्ट वाटेल.
\end{explain}

\begin{thm}
\label{cmp:thy:smn:thm:s-m-n}
  नैसर्गिक संख्यांच्या प्रत्येक $n$ आणि~$m$ जोडीसाठी असे आदिम
  पुनरावर्ती फलन~$s^m_n$ असते की प्रत्येक क्रमिका
  $e$, $a_0$, \dots, $a_{m-1}$, $y_0$ ,\dots, $y_{n-1}$ साठी
  \[  \cfind{s^m_n(e, a_0, \dots, a_{m-1})}[n](y_0, \dots, y_{n-1}) \simeq
  \cfind{e}[m+n](a_0, \dots, a_{m-1}, y_0, \dots, y_{n-1}).
\]
\end{thm}

\begin{explain}
$s^m_n$ हे \emph{कार्यक्रमांवर} कार्य करते असा विचार करणे उपयुक्त आहे.
म्हणजे, $s^m_n$ हे $(m+n)$-स्थानी फलनाचा कार्यक्रम~$e$ आणि निश्चित
आदाने $a_0$, \dots, $a_{m-1}$ घेते; उरलेल्या आदानांच्या $n$-स्थानी
फलनासाठी ते कार्यक्रम
$s^m_n(e, a_0, \dots, a_{m-1})$ परत करते.


\end{explain}











\section{सार्वत्रिक आंशिक संगणनक्षम फलन}\label{cmp:thy:uni:sec}

\begin{thm}
\label{cmp:thy:uni:thm:univ-comp}
एक सार्वत्रिक आंशिक संगणनक्षम फलन~$\fn{Un}(e,x)$ असते. दुसऱ्या
शब्दांत, पुढील गुणधर्म असलेले फलन $\fn{Un}(e,x)$ असते:
\begin{enumerate}
\item $\fn{Un}(e,x)$ आंशिक संगणनक्षम आहे.
\item $f(x)$ हे कोणतेही आंशिक संगणनक्षम फलन असेल, तर अशी नैसर्गिक
संख्या $e$ असते की प्रत्येक~$x$ साठी $f(x) \simeq \fn{Un}(e,x)$.
\end{enumerate}
\end{thm}

\begin{proof}
क्लिनी यांच्या प्रमाण रूप प्रमेयातील (\ref{cmp:thy:nfm:thm:normal-form})
$U$ आणि $T$ घेऊन $\fn{Un}(e,x) \simeq U(\umin{s}{T(e,x,s)})$
अशी व्याख्या करा.
\end{proof}

\begin{explain}
आंशिक संगणनक्षम फलनांचे प्रभावी प्रगणन आहे, हे अचूकपणे सांगण्याचा
हा मार्ग आहे. कल्पना अशी की $f_e(x) = \fn{Un}(e,x)$ ने परिभाषित
केलेल्या फलनाला $f_e$ म्हटले, तर $f_0$, $f_1$, $f_2$, \dots या
क्रमिकेत सर्व आंशिक संगणनक्षम फलने येतात; शिवाय $f_e(x)$ चे
संगणन $e$ आणि~$x$ या दोहोंवर ``एकसमानपणे'' करता येते. सोपेपणासाठी
आपण एकस्थानी फलनांसाठी सार्वत्रिक असलेले द्विस्थानी फलन वापरत आहोत.
परंतु संख्यांच्या क्रमिकांचे सांकेतिकरण करून हे अधिक आदानांसाठी सहज
सामान्य करता येते. उदाहरणार्थ, $f(x,y,z)$ हे $3$-स्थानी आंशिक
पुनरावर्ती फलन असेल, तर $g(x) \simeq f((x)_0, (x)_1, (x)_2)$
हे एकस्थानी आंशिक पुनरावर्ती फलन असते.

\end{explain}











\section{सार्वत्रिक संगणनक्षम फलन नसते}\label{cmp:thy:nou:sec}

आंशिक संगणनक्षम फलनांसाठी सार्वत्रिक असे आंशिक संगणनक्षम फलन असले,
तरी सर्वत्र परिभाषित संगणनक्षम फलनांसाठी सार्वत्रिक असे सर्वत्र
परिभाषित संगणनक्षम फलन नसते.


\begin{thm}
\label{cmp:thy:nou:thm:no-univ}
सार्वत्रिक संगणनक्षम फलन नसते. दुसऱ्या शब्दांत, असे कोणतेही
फलन $\fn{Un}'(k, x)$ घ्या की $f(x)$ हे सर्वत्र परिभाषित संगणनक्षम
फलन असल्यास, अशी नैसर्गिक संख्या~$k$ असते की प्रत्येक~$x$ साठी
$f(x) = \fn{Un}'(k,x)$; तर असे फलन संगणनक्षम नसते.
\end{thm}

\begin{proof}
सिद्धता साध्या विकर्णीकरणाने मिळते: $\fn{Un}'(k,x)$ हे सर्वत्र
परिभाषित आणि संगणनक्षम असेल, तर
\[d(x) = \fn{Un}'(x, x) + 1
\]
हेही सर्वत्र परिभाषित आणि संगणनक्षम असेल. परंतु व्याख्येवरून
$d(k)$ हे $\fn{Un}'(k,k)$ च्या समान नाही. म्हणून प्रत्येक $k$
साठी $d(x)$ आणि~$\fn{Un}'(k, x)$ यांची मूल्ये किमान एका~$x$
साठी, म्हणजे $x = k$ साठी, वेगळी असतात.
\end{proof}

\begin{explain}
वरील \ref{cmp:thy:uni:thm:univ-comp} दाखवते की सार्वत्रिक फलन आंशिक
असण्याची परवानगी दिल्यास हे विकर्णीकरण टाळता येते. म्हणजे,
$\fn{Un}$ हे सर्वत्र परिभाषित संगणनक्षम फलनांसाठी सार्वत्रिक आहे;
पण ते स्वतः सर्वत्र परिभाषित नाही. आंशिक फलनांच्या बाबतीत
विकर्णीकरणाचा युक्तिवाद चालत नाही.
\end{explain}

\begin{prob}
  \ref{cmp:thy:nou:thm:no-univ} च्या सिद्धतेतील विकर्णीकरणाचा युक्तिवाद आंशिक
  फलनांच्या बाबतीत का चालत नाही हे समजण्यासाठी
  $f(x) \simeq \fn{Un}(x,x)+1$ हे फलन विचारात घ्या. ते आंशिक
  संगणनक्षम आहे का? असल्यास, त्याला निर्देशांक~$e$ आहे, म्हणजे
  $f(x) \simeq \fn{Un}(e,x)$. मग~$f(e)$ विषयी काय म्हणता येईल?
\end{prob}












\section{थांबण्याची समस्या}\label{cmp:thy:hlt:sec}

रचनेवरूनच, सार्वत्रिक आंशिक संगणनक्षम फलन~$\fn{Un}(e,x)$ तेव्हा आणि
केवळ तेव्हाच परिभाषित असते, जेव्हा~$e$ ने संकेतबद्ध केलेल्या फलनाच्या आदान~$x$
वरील संगणनातून मूल्य मिळते. असे होते की नाही हे आपण ठरवू शकतो का,
हा नैसर्गिक प्रश्न आहे. प्रत्यक्षात ते ठरवता येत नाही. ट्यूरिंग
यंत्रांच्या संगणन-प्रतिमानात याचा अर्थ, दिलेले ट्यूरिंग यंत्र दिलेल्या
आदानावर थांबते का हे संगणनाने ठरवता येत नाही. म्हणून पुढील प्रमेयाला
``थांबण्याच्या समस्येची अनिर्णेयता'' म्हणतात. खाली आपण दोन सिद्धता
देऊ. पहिली आपल्या आधीच्या चर्चेचा धागा पुढे नेते, तर दुसरी अधिक थेट आहे.

\begin{thm}
\label{cmp:thy:hlt:thm:halting-problem}
पुढीलप्रमाणे असू द्या:
\[h(e, x)  =
\begin{cases}
1 & \text{जर\/ $\fn{Un}(e, x)$ परिभाषित असेल तर} \\
0 & \text{अन्यथा.}
\end{cases}
\]
मग $h$ संगणनक्षम नाही.
\end{thm}

\begin{proof}
$h$ संगणनक्षम आहे असे समजा. मग सार्वत्रिक संगणनक्षम फलन परिभाषित
करता येईल, हे दाखवू. पुढीलप्रमाणे व्याख्या करा:
\[\fn{Un'}(e,x) =
\begin{cases}
\fn{Un}(e,x) & \text{जर $h(e,x) = 1$ असेल तर} \\
0 & \text{अन्यथा.}
\end{cases}
\]
आता $\fn{Un'}(e, x)$ हे सर्वत्र परिभाषित फलन आहे आणि $h$
संगणनक्षम असेल तर तेही संगणनक्षम आहे. उदाहरणार्थ, आदिम पुनरावर्तन
वापरून $g$ ची व्याख्या पुढीलप्रमाणे करता येते:
\begin{align*}
g(0, e, x) & \simeq 0 \\
g(y+1, e, x) & \simeq \fn{Un}(e,x);
\end{align*}
मग
\[\fn{Un'}(e,x) \simeq g(h(e,x),e,x).
\]
$\fn{Un'}(e,x)$ हे $\fn{Un}(e,x)$ परिभाषित असेल तेथे त्याच्याशी
सहमत असल्याने, $\fn{Un'}$ हे सर्वत्र परिभाषित असलेल्या आंशिक
संगणनक्षम फलनांसाठी सार्वत्रिक आहे. पण यामुळे
\ref{cmp:thy:nou:thm:no-univ} शी विरोध होतो.
\end{proof}

\begin{proof}
$h(e,x)$ संगणनक्षम आहे असे समजा. फलन $g$ ची पुढीलप्रमाणे व्याख्या करा:
\[g(x) =
\begin{cases}
  0                & \text{जर $h(x,x) = 0$ असेल तर} \\
  \fundefined & \text{अन्यथा.}
\end{cases}
\]
फलन $g$ आंशिक संगणनक्षम आहे. उदाहरणार्थ, त्याची व्याख्या
$\umin{y}{h(x,x) = 0}$ अशी करता येते. म्हणून एखाद्या~$e$ साठी
प्रत्येक~$x$ बाबत $g(x) \simeq \fn{Un}(e,
x)$ असते. $g$ हे $e$ येथे परिभाषित आहे का? असल्यास, $g$ च्या
व्याख्येवरून $h(e,e) = 0$ ($h$ परिभाषित असेल तेव्हाच ते~$0$ मूल्य
घेऊ शकते). $h$ च्या व्याख्येनुसार, याचा अर्थ $\fn{Un}(e, e)$
अपरिभाषित आहे. पण प्रत्येक~$x$ साठी $g(x) \simeq \fn{Un}(e, x)$
या गृहितकावरून $g(e)$ अपरिभाषित ठरते—विरोध. दुसरीकडे, $g(e)$
अपरिभाषित असेल, तर $h(e,e) \neq 0$, म्हणून $h(e,e) = 1$.
त्यामुळे $\fn{Un}(e, e)$ परिभाषित आहे. पण
$g(x) \simeq \fn{Un}(e, x)$ असल्यामुळे $g(e)$ देखील परिभाषित
असेल. पुन्हा विरोध मिळतो.
\end{proof}













\section{रसेलच्या विरोधापत्तीशी तुलना}\label{cmp:thy:rus:sec}

या विभागातील युक्तिवादांची रसेलच्या विरोधापत्तीशी तुलना केल्यास
त्यांतील साम्य आणि भेद स्पष्ट होतात:
\begin{enumerate}
\item रसेलची विरोधापत्ती: $S = \Setabs{x}{x \notin x}$ असे मानू.
  मग $S
  \in S$ तेव्हा आणि केवळ तेव्हाच, जेव्हा $S \notin S$; हा विरोध आहे.

  \emph{निष्कर्ष:} असा संच~$S$ अस्तित्वात नाही. ``सर्व संचांचा संच''
  अस्तित्वात आहे हे गृहितक संचसिद्धांताच्या इतर स्वयंसिद्धांशी विसंगत आहे.


\item रसेलच्या विरोधापत्तीचे एक रूपांतर: सर्व फलनांच्या संचापासून
  $\{ 0, 1 \}$ मध्ये जाणारे ``फलन'' $F$ पुढीलप्रमाणे परिभाषित करू:
  \[  F(f) =
  \begin{cases}
    1 & \text{जर $f$ हे $f$ च्या परिभाषाक्षेत्रात असेल आणि $f(f) = 0$ असेल तर} \\
    0 & \text{अन्यथा}
  \end{cases}
  \]
  तशाच युक्तिवादावरून $F(F) = 0$ तेव्हा आणि केवळ तेव्हाच, जेव्हा
  $F(F) = 1$; हा विरोध आहे.

  \emph{निष्कर्ष:} $F$ हे फलन नाही. ``सर्व फलनांचा संच'' इतका
  मोठा आहे की तो फलनाचे परिभाषाक्षेत्र असू शकत नाही.

\item विकर्णीकरणाचा युक्तिवाद: $f_0$, $f_1$, \dots हे आंशिक
  संगणनक्षम फलनांचे प्रगणन असू द्या आणि $G \colon \Nat \to
  \{ 0, 1 \}$ ची व्याख्या पुढीलप्रमाणे करू:
  \[  G(x) =
  \begin{cases}
    1 & \text{जर $f_x(x)\downarrow = 0$ असेल तर} \\
    0 & \text{अन्यथा}
  \end{cases}
  \]
  $G$ संगणनक्षम असेल, तर एखाद्या $k$ साठी ते फलन $f_k$
  असेल. पण मग $G(k) = 1$ तेव्हा आणि केवळ तेव्हाच, जेव्हा $G(k) = 0$;
  हा विरोध आहे.

  \emph{निष्कर्ष:} $G$ संगणनक्षम नाही. संचसिद्धांताच्या
  स्वयंसिद्धांनुसार $G$ हे तरीही फलन आहे, हे लक्षात घ्या;
  येथे विरोधापत्ती नाही, तर फरक स्पष्ट झाला आहे.
\end{enumerate}

आंशिक फलने, संगणनक्षम फलने, आंशिक संगणनक्षम फलने इत्यादींबद्दलची
चर्चा गोंधळात टाकू शकते. $\Nat$ पासून $\Nat$ मध्ये जाणाऱ्या सर्व
आंशिक फलनांचा संच हा वस्तूंचा मोठा संग्रह आहे. त्यांपैकी काही
सर्वत्र परिभाषित आहेत, काही आंशिक संगणनक्षम आहेत, काही सर्वत्र
परिभाषित आणि संगणनक्षम अशी दोन्ही आहेत, तर काही यांपैकी कोणतीच नाहीत.
आपण नुसते ``फलन'' म्हटले, तर सामान्यतः सर्वत्र परिभाषित फलन अभिप्रेत
असते, हे लक्षात ठेवा. म्हणून पुढील प्रकार मिळतात:
\begin{enumerate}
\item संगणनक्षम फलने
\item सर्वत्र परिभाषित नसलेली आंशिक संगणनक्षम फलने
\item संगणनक्षम नसलेली फलने
\item सर्वत्र परिभाषितही नसलेली आणि संगणनक्षमही नसलेली आंशिक फलने
\end{enumerate}
हा फरक समजण्यासाठी $\Nat$ पासून $\Nat$ मध्ये जाणारी सर्व आंशिक
फलने दर्शवणारा मोठा चौकोन काढा. त्यात सर्वत्र परिभाषित फलने
आणि आंशिक संगणनक्षम फलने दर्शवणारी, एकमेकांवर काही प्रमाणात
येणारी दोन क्षेत्रे दाखवा. अशा आकृतीत मिळणाऱ्या प्रत्येक
क्षेत्रातील एखाद्या वस्तूचे वर्णन करता येते का, हे पाहणे चांगला सराव आहे.











\section{संगणनक्षम संच}\label{cmp:thy:cps:sec}

संगणनक्षमतेची कल्पना संगणनक्षम फलनांपासून संगणनक्षम संचांपर्यंत
विस्तारता येते:

\begin{defn}
  $S$ हा नैसर्गिक संख्यांचा संच असू द्या. $S$ \emph{संगणनक्षम}
  असतो तेव्हा आणि केवळ तेव्हाच, जेव्हा त्याचे जातिबोधक फलन~$\Char{S}$
  संगणनक्षम असते. दुसऱ्या शब्दांत, $S$~संगणनक्षम असतो तेव्हा आणि केवळ
  तेव्हाच, जेव्हा पुढील फलन संगणनक्षम असते:
\[\Char{S}(x) =
\begin{cases}
1 & \text{जर $x \in S$ असेल तर} \\
0 & \text{अन्यथा}
\end{cases}
\]
त्याचप्रमाणे, संबंध $R(x_0, \dots, x_{k-1})$ संगणनक्षम असतो तेव्हा आणि
केवळ तेव्हाच, जेव्हा त्याचे जातिबोधक फलन संगणनक्षम असते.

संगणनक्षम संच आणि संबंध यांना \emph{निर्णेय} असेही म्हणतात.
\end{defn}

\begin{explain}
आता संगणनक्षमतेच्या अनेक कल्पना आहेत: आंशिक फलनांसाठी, फलनांसाठी
आणि संचांसाठी. त्यांची गल्लत करू नका!{}
\end{explain}











\section{संगणनक्षमपणे प्रगणनीय संच}\label{cmp:thy:ces:sec}

\begin{defn}
एखादा संच रिकामा असेल किंवा एखाद्या संगणनक्षम फलनाची व्याप्ती असेल,
तर तो \emph{संगणनक्षमपणे प्रगणनीय} असतो.
\end{defn}

\begin{history}
संगणनक्षमपणे प्रगणनीय संचांना \emph{पुनरावर्तीपणे प्रगणनीय}
असेही म्हणतात. ही मूळ संज्ञा आहे. आज ही दोन्ही नावे, तसेच
``c.e.'' आणि ``r.e.'' ही संक्षिप्त रूपे, सर्वसाधारणपणे वापरली जातात.
\end{history}

\begin{explain}
या व्याख्येचा अर्थ काय आणि ही संज्ञा योग्य का आहे, याचा विचार करा.
कल्पना अशी की $S$ हा संगणनक्षम फलन~$f$ ची व्याप्ती असेल, तर
\[S = \{ f(0), f(1), f(2), \dots \},
\]
म्हणून $f$ हे~$S$ च्या घटकांचे ``प्रगणन'' करते असे मानता येते.
व्याख्येनुसार $f$ वाढते फलन असणे आवश्यक नाही, म्हणजे प्रगणन
चढत्या क्रमाने असणे आवश्यक नाही. किंबहुना, $f$ एकास-एक असणेही
आवश्यक नाही; म्हणजे $f(0)$, $f(1)$, $f(2)$, \dots{} या~$S$
च्या प्रगणनात पुनरुक्ती चालते. उदाहरणार्थ, स्थिर फलन $f(x) = 0$
हे संच $\{ 0
\}$ चे प्रगणन करते.
\end{explain}

प्रत्येक संगणनक्षम संच संगणनक्षमपणे प्रगणनीय असतो. हे पाहण्यासाठी,
$S$~संगणनक्षम आहे असे समजा. $S$ रिकामा असेल, तर व्याख्येवरून
तो संगणनक्षमपणे प्रगणनीय आहे. अन्यथा $a$ हा $S$ मधील कोणताही
घटक घ्या. $f$ ची पुढीलप्रमाणे व्याख्या करा:
\[f(x) =
\begin{cases}
x & \text{जर $\Char{S}(x) = 1$ असेल तर} \\
a & \text{अन्यथा.}
\end{cases}
\]
मग $f$ हे संगणनक्षम फलन आहे आणि $S$ ही~$f$ ची व्याप्ती आहे.












\section{संगणनक्षमपणे
  प्रगणनीय संचांच्या समतुल्य व्याख्या}\label{cmp:thy:eqc:sec}


संगणनक्षमपणे प्रगणनीय असण्याचा अर्थ सांगणारी पुढील काही महत्त्वाची
समतुल्य विधाने आहेत.

\begin{thm}
\label{cmp:thy:eqc:thm:ce-equiv}
$S$ हा नैसर्गिक संख्यांचा संच असू द्या. मग पुढील विधाने समतुल्य आहेत:
\begin{enumerate}
\item\label{cmp:thy:eqc:case:ce} $S$ संगणनक्षमपणे प्रगणनीय आहे.
\item\label{cmp:thy:eqc:case:ran-pc} $S$ ही एखाद्या \emph{आंशिक}
  संगणनक्षम फलनाची व्याप्ती आहे.
\item\label{cmp:thy:eqc:case:ran-prim} $S$ रिकामा आहे किंवा एखाद्या आदिम
  पुनरावर्ती फलनाची व्याप्ती आहे.
\item\label{cmp:thy:eqc:case:ce-domain} $S$ हे एखाद्या आंशिक संगणनक्षम
  फलनाचे \emph{परिभाषाक्षेत्र} आहे.
\end{enumerate}
\end{thm}

\begin{explain}
पहिली तीन विधाने सांगतात की कोणत्याही रिकाम्या नसलेल्या
संगणनक्षमपणे प्रगणनीय संचाचे प्रगणन संगणनक्षम फलनाने, आंशिक
संगणनक्षम फलनाने किंवा आदिम पुनरावर्ती फलनाने करता येते.
चौथे विधान सांगते की $S$ संगणनक्षमपणे प्रगणनीय असेल, तर
एखाद्या निर्देशांक~$e$ साठी
\[S = \Setabs{x}{\cfind{e}(x) \fdefined}.
\]
दुसऱ्या शब्दांत, $S$ हा अशा आदानांचा संच आहे ज्यांवर
$\cfind{e}$ चे संगणन थांबते. म्हणून संगणनक्षमपणे प्रगणनीय
संचांना कधी कधी \emph{अर्धनिर्णेय} म्हणतात: संख्या त्या संचात
असेल, तर शेवटी ``होय'' असे उत्तर मिळते; पण ती संचात नसेल,
तर ``नाही'' असे उत्तर कधीच मिळत नाही!{}
\end{explain}

\begin{proof}
प्रत्येक आदिम पुनरावर्ती फलन संगणनक्षम असते आणि प्रत्येक
संगणनक्षम फलन आंशिक संगणनक्षमही असते. म्हणून
\ref{cmp:thy:eqc:case:ran-prim} वरून \ref{cmp:thy:eqc:case:ce} मिळते, आणि
\ref{cmp:thy:eqc:case:ce} वरून~\ref{cmp:thy:eqc:case:ran-pc} मिळते.
($S$ रिकामा असेल, तर $S$ कुठेही परिभाषित नसलेल्या आंशिक
संगणनक्षम फलनाची व्याप्ती असतो, हे लक्षात घ्या.)
\ref{cmp:thy:eqc:case:ran-pc} वरून \ref{cmp:thy:eqc:case:ran-prim} मिळते हे
दाखवले, तर पहिली तीन विधाने समतुल्य ठरतील.

म्हणून $S$ ही आंशिक संगणनक्षम फलन $\cfind{e}$ ची व्याप्ती आहे
असे समजा. $S$ रिकामा असेल, तर काम झाले. अन्यथा $a$ हा~$S$
मधील कोणताही घटक घ्या. क्लिनी यांच्या प्रमाण रूप प्रमेयावरून
पुढीलप्रमाणे लिहिता येते:
\[\cfind{e}(x) = U(\umin{s}{T(e, x, s)}).
\]
विशेषतः, $\cfind{e}(x) \fdefined$ असून त्याचे मूल्य $= y$ असते तेव्हा
आणि केवळ तेव्हाच, जेव्हा असा $s$ असतो की $T(e, x, s)$ आणि $U(s) = y$.
$f(z)$ ची व्याख्या पुढीलप्रमाणे करा:
\[f(z) = \begin{cases}
  U((z)_1) & \text{जर $T(e, (z)_0, (z)_1)$ असेल तर} \\
  a        & \text{अन्यथा.}
\end{cases}
\]
मग $f$ आदिम पुनरावर्ती आहे, कारण $T$ आणि $U$
तशी आहेत.

$S$ ही~$f$ ची व्याप्ती आहे हे दाखवायचे आहे; म्हणजे प्रत्येक
नैसर्गिक संख्या~$y$ साठी $y \in S$ तेव्हा आणि केवळ तेव्हाच, जेव्हा
ती~$f$ च्या व्याप्तीत असते. प्रथम $y \in S$ असे समजा. मग $y$ ही
$\cfind{e}$ च्या व्याप्तीत आहे. म्हणून काही $x$ आणि~$s$
साठी $T(e,x,s)$ लागू होतो आणि $U(s) = y$. त्यामुळे
$y = f(\tuple{x,s})$. उलट, $y$ ही~$f$ च्या व्याप्तीत
असेल, तर एकतर $y = a$ असेल, किंवा एखाद्या~$z$ साठी
$T(e,(z)_0,(z)_1)$ लागू होईल आणि $U((z)_1) = y$ असेल.
दुसऱ्या प्रकरणात $\cfind{e}((z)_0) \fdefined =
y$ असल्याने, दोन्ही प्रकारे $y$ ही~$S$ मध्येच असते.


($\cfind{e}(x) \fdefined = y$ या चिन्हांकनाचा अर्थ
``$\cfind{e}(x)$ परिभाषित आहे आणि $y$ च्या समान आहे'' असा आहे.
त्याऐवजी $\cfind{e}(x) =
y$ असेही लिहिता येईल; पण आपण आंशिक फलनाशी व्यवहार करत आहोत,
याची आठवण करून देण्यासाठी अतिरिक्त बाण कधी कधी उपयुक्त असतो.)

\ref{cmp:thy:eqc:thm:ce-equiv} ची सिद्धता पूर्ण करण्यासाठी
\ref{cmp:thy:eqc:case:ce} आणि~\ref{cmp:thy:eqc:case:ce-domain} समतुल्य आहेत
हे दाखवणे पुरेसे आहे. प्रथम \ref{cmp:thy:eqc:case:ce} वरून
\ref{cmp:thy:eqc:case:ce-domain} मिळते हे पाहू. संच रिकामा असेल,
तर कुठेही परिभाषित नसलेल्या आंशिक संगणनक्षम फलनाचे तो
परिभाषाक्षेत्र आहे. आता तो रिकामा नाही असे समजा. मग
$S$ ही संगणनक्षम फलन~$f$ ची व्याप्ती आहे, म्हणजे
\[S = \Setabs{y}{\text{अशा काही $x$ साठी, } f(x) = y}.
\]
पुढीलप्रमाणे घ्या:
\[g(y) = \umin{x}{(f(x) = y)}.
\]
मग $g$ आंशिक संगणनक्षम आहे आणि $g(y)$ तेव्हा आणि केवळ तेव्हाच
परिभाषित असते, जेव्हा काही~$x$ साठी $f(x) = y$ असते. दुसऱ्या शब्दांत, $g$
चे परिभाषाक्षेत्र ही~$f$ ची व्याप्ती आहे.



शेवटी \ref{cmp:thy:eqc:case:ce-domain} वरून~\ref{cmp:thy:eqc:case:ce} मिळते
हे दाखवू. $S$ हे आंशिक संगणनक्षम फलन~$\cfind{e}$ चे
परिभाषाक्षेत्र आहे असे समजा, म्हणजे
\[S = \Setabs{x}{\cfind{e}(x) \fdefined}.
\]
$S$ रिकामा असेल, तर काम झाले; अन्यथा $a$ हा~$S$
मधील कोणताही घटक घ्या. $f$ ची पुढीलप्रमाणे व्याख्या करा:
\[f(z) = \begin{cases}
(z)_0 & \text{जर $T(e,(z)_0,(z)_1)$ असेल तर} \\
a & \text{अन्यथा.}
\end{cases}
\]
मग वरच्याप्रमाणे, संख्या $x$ ही~$f$ च्या व्याप्तीत असते तेव्हा आणि
केवळ तेव्हाच, जेव्हा $\cfind{e}(x) \fdefined$; आणि ही अट तेव्हा आणि
केवळ तेव्हाच खरी असते, जेव्हा $x \in S$.

\end{proof}

\ref{cmp:thy:eqc:thm:ce-equiv} मधील विधान~\ref{cmp:thy:eqc:case:ce-domain}
संगणनक्षमपणे प्रगणनीय संचांचे प्रगणन करण्याचा सोयीचा मार्ग देते:
प्रत्येक~$e$ साठी $W_e$ हे $\cfind{e}$ चे परिभाषाक्षेत्र दर्शवू,
म्हणजे
\[W_e = \Setabs{x}{\cfind{e}(x) \fdefined}.
\]
मग $A$ हा कोणताही संगणनक्षमपणे प्रगणनीय संच असेल, तर
एखाद्या~$e$ साठी $A = W_e$.

पुढील प्रमेय संगणनक्षमपणे प्रगणनीय संचांचे आणखी एक वैशिष्ट्य
सांगते.

\begin{thm}
\label{cmp:thy:eqc:thm:exists-char}
$S$ हा संच संगणनक्षमपणे प्रगणनीय असतो तेव्हा आणि केवळ तेव्हाच, जेव्हा
असा संगणनक्षम संबंध $R(x,y)$ असतो की
\[S = \Setabs{ x }{ \lexists[y][R(x,y)] }.
\]
\end{thm}

\begin{proof}
प्रथम $S$ संगणनक्षमपणे प्रगणनीय आहे असे समजा. मग काही
$e$ साठी $S = W_e$. या~$e$ साठी $S$ असे लिहिता येते:
\[S = \Setabs{ x }{ \lexists[y][T(e, x, y)] }.
\]
उलट, $S = \Setabs{ x }{
  \lexists[y][R(x, y)] }$ असे समजा. $f$ ची पुढीलप्रमाणे व्याख्या करा:
\[f(x) \simeq \umin{y}{R(x, y)}.
\]
मग $f$ आंशिक संगणनक्षम आहे आणि $S$ हे~$f$ चे परिभाषाक्षेत्र आहे.
\end{proof}












\section{संगणनक्षम नसलेले संच अस्तित्वात आहेत}\label{cmp:thy:ncp:sec}


प्रत्येक संगणनक्षम संच संगणनक्षमपणे प्रगणनीय असतो, हे आपण वर पाहिले.
याचे उलटही खरे आहे का? पुढील प्रमेय दाखवते की सर्वसाधारणपणे तसे नाही.

\begin{thm}\label{cmp:thy:ncp:thm:K-0}
$K_0$ हा $\Setabs{\tuple{e, x}}{\cfind{e}(x) \fdefined}$ हा
संच असू द्या. मग $K_0$ संगणनक्षमपणे प्रगणनीय आहे, पण संगणनक्षम नाही.
\end{thm}

\begin{proof}
$K_0$ संगणनक्षमपणे प्रगणनीय आहे हे पाहण्यासाठी, तो पुढीलप्रमाणे
परिभाषित फलन~$f$ चे परिभाषाक्षेत्र आहे हे लक्षात घ्या:
\[f(z) = \umin{y}{(z = \tuple{(z)_0,(z)_1} \land T((z)_0, (z)_1, y))}.
\]
कारण $\cfind{e}(x)$ परिभाषित असेल, तर $f(\tuple{e, x})$
थांबणारी संगणनक्रमिका शोधते; $\cfind{e}(x)$ अपरिभाषित असेल,
तर $f(\tuple{e, x})$ देखील अपरिभाषित असते; आणि $z$ हा
जोडीचा संकेतांकसुद्धा नसेल, तर $f(z)$ देखील अपरिभाषित असते. केवळ लांबी
दोन येते एवढी कसोटी पुरेशी नाही: 42 या अवैध संकेतांकासाठी लांबीचे फलन
दोन देते, पण दोन्ही काढलेल्या घटकांना पुन्हा संकेतबद्ध केल्यास 6 मिळते.
म्हणून येथे जोडीचे अचूक पुनःसंकेतबद्धीकरण तपासले आहे.


$K_0$ संगणनक्षम नाही, ही बाब म्हणजे थांबण्याच्या समस्येचीच
अनिर्णेयता, \ref{cmp:thy:hlt:thm:halting-problem}.
\end{proof}

संच $K_0$ म्हणजे $\cfind{e}(x) \fdefined$ असा गुणधर्म असलेल्या
$\tuple{e,x}$ जोड्यांचा संच. म्हणजे $\tuple{e,x} \in K_0$
तेव्हा आणि केवळ तेव्हाच, जेव्हा $\cfind{e}$ हे आदान~$x$ वर परिभाषित असते
(संगणन थांबते). म्हणून त्याला ``थांबणारा संच'' असेही म्हणतात.
$K = \Setabs{e}{\cfind{e}(e) \fdefined}$ हा संच
``स्वतःच्या निर्देशांकावर थांबणारा संच'' आहे. तो अनिर्णेय संचाचे
मानक उदाहरण म्हणून वारंवार वापरला जातो.

\begin{thm}\label{cmp:thy:ncp:thm:K}
स्वतःच्या निर्देशांकावर थांबणारा संच $K = \Setabs{e}{\cfind{e}(e) \fdefined}$
हा संगणनक्षमपणे प्रगणनीय आहे, पण निर्णेय नाही.
\end{thm}

\begin{proof}
  $K$ निर्णेय आहे, म्हणजे त्याचे जातिबोधक फलन
  $\Char{K}$ संगणनक्षम आहे, असे समजा. पुढीलप्रमाणे घ्या:
  \[d(e) = \begin{cases}
  1 & \text{जर\/ $\Char{K}(e) = 0$ असेल तर}\\
  \fundefined & \text{अन्यथा.}
  \end{cases}
  \]
  $k$ हा~$d$ चा निर्देशांक असू द्या, म्हणजे $d \simeq \cfind{k}$.
  मग $d(k)
  \simeq \cfind{k}(k)$. पण $d(k)
  \fdefined$ तेव्हा आणि केवळ तेव्हाच, जेव्हा $\cfind{k}(k) \fundefined$, हे~$d$
  च्या व्याख्येवरून मिळते; त्यामुळे विरोध होतो.

  $K$ हे $f(x) = \umin{y}{T(x,x,y)}$ चे परिभाषाक्षेत्र आहे,
  म्हणून ते संगणनक्षमपणे प्रगणनीय आहे.
\end{proof}












\section{संगणनक्षमपणे
  प्रगणनीय संच संयोग आणि छेद यांखाली बंद असतात}\label{cmp:thy:clo:sec}

पुढील प्रमेय संगणनक्षमपणे प्रगणनीय संचांच्या काही संवरण गुणधर्मांबद्दल आहे.

\begin{thm}
$A$ आणि $B$ हे संगणनक्षमपणे प्रगणनीय असतील, तर
$A \cap B$ आणि $A \cup B$ देखील तसेच असतात.
\end{thm}

\begin{proof}
\ref{cmp:thy:eqc:thm:ce-equiv} मुळे संगणनक्षमपणे प्रगणनीय संचांची
वेगवेगळी वैशिष्ट्ये वापरता येतात. त्याचे उदाहरण म्हणून काही
वेगवेगळ्या सिद्धता देऊ.

पहिल्या दोन सिद्धतांत दोन्ही संच रिकामे नाहीत असे धरा. एखादा
रिकामा असेल, तर त्यांचा छेद रिकामा आणि संयोग दुसरा संच असतो;
निष्कर्ष त्या वेळी थेट मिळतो.

पहिल्या सिद्धतेसाठी $A$ चे प्रगणन संगणनक्षम फलन~$f$ करते
आणि $B$ चे प्रगणन संगणनक्षम फलन~$g$ करते असे समजा.
पुढीलप्रमाणे घ्या:
\begin{align*}
h(x) & = \umin{y}{(f(y) = x \lor g(y) = x)} \text{ आणि}\\
j(x) & = \umin{y}{(f((y)_0) = x \land g((y)_1) = x)}.
\end{align*}
मग $A \cup B$ हे $h$ चे परिभाषाक्षेत्र आहे आणि
$A \cap B$ हे~$j$ चे परिभाषाक्षेत्र आहे.

\begin{explain}
संगणनाच्या भाषेत येथे पुढील गोष्ट घडते: $A$ आणि $B$ यांचे
प्रगणन करणाऱ्या प्रक्रिया दिल्या असता, घटक $x$ हा $A \cup B$
मध्ये आहे का हे एका तरी प्रगणनात~$x$ शोधून अर्धनिर्णीत करता
येते. तसेच $x$ हा $A \cap B$ मध्ये आहे का हे दोन्ही
प्रगणनांत एकाच वेळी~$x$ शोधून अर्धनिर्णीत करता येते.
\end{explain}

दुसऱ्या सिद्धतेसाठी पुन्हा $A$ चे प्रगणन~$f$ करते आणि
$B$ चे प्रगणन~$g$ करते असे समजा. पुढीलप्रमाणे घ्या:
\[k(x) = \begin{cases}
f(x/2) & \text{जर $x$ सम असेल तर} \\
g((x-1)/2) & \text{जर $x$ विषम असेल तर.}
\end{cases}
\]
मग $k$ हे $A \cup B$ चे प्रगणन करते. कल्पना एवढीच की
$k$ हे $f$ आणि~$g$ यांच्या प्रगणनांचा आलटूनपालटून वापर करते.
$A \cap
B$ चे प्रगणन अधिक अवघड आहे. $A \cap B$ रिकामा असेल,
तर तो अर्थातच संगणनक्षमपणे प्रगणनीय आहे. अन्यथा $c$ हा
$A \cap B$ मधील कोणताही घटक घ्या आणि $l$ ची व्याख्या
पुढीलप्रमाणे करा:
\[l(x) = \begin{cases}
f((x)_0) & \text{जर $f((x)_0) = g((x)_1)$ असेल तर} \\
c & \text{अन्यथा.}
\end{cases}
\]
संगणनाच्या भाषेत, $l$ हे $f$ आणि $g$ यांच्या प्रगणनांतील
घटकांच्या जोड्यांमधून जाते आणि सापडलेली प्रत्येक जुळणी परत
करते; अन्यथा ते~$c$ परत करून प्रगणनात प्रगती करत नाही.

शेवटच्या सिद्धतेसाठी $A$ हे आंशिक फलन $m(x)$ चे
\emph{परिभाषाक्षेत्र} आणि $B$ हे आंशिक फलन $n(x)$ चे
परिभाषाक्षेत्र आहे असे समजा. मग $A \cap B$ हे आंशिक फलन
$m(x) +
n(x)$ चे परिभाषाक्षेत्र आहे.

\begin{explain}
संगणनाच्या भाषेत, $A$ हा ज्यांवर $m$ थांबते त्या
मूल्यांचा संच असेल आणि $B$ हा ज्यांवर $n$ थांबते त्या
मूल्यांचा संच असेल, तर $A \cap B$ हा दोन्ही प्रक्रिया
थांबतात अशा मूल्यांचा संच आहे.
\end{explain}

$A \cup B$ हा थांबणाऱ्या आदानांचा संच म्हणून व्यक्त करणे
अधिक कठीण आहे, कारण $m$ आणि $n$ यांचे समांतर अनुकरण
करावे लागते. $d$ हा $m$ चा निर्देशांक आणि $e$ हा $n$
चा निर्देशांक असू द्या; दुसऱ्या शब्दांत, $m =
\cfind{d}$ आणि $n = \cfind{e}$. मग $A \cup B$ हे
पुढील फलनाचे परिभाषाक्षेत्र आहे:
\[p(x) = \umin{y}{(T(d,x,y) \lor T(e,x,y))}.
\]
\begin{explain}
संगणनाच्या भाषेत, आदान $x$ वरील $p$ हे $m$ किंवा
$n$ यांपैकी एखाद्याचे थांबणारे संगणन शोधते आणि त्यांपैकी
एक सापडल्यास थांबते.
\end{explain}
\end{proof}











\section{पूरक संचाखाली संगणनक्षमपणे प्रगणनीय संच बंद नसतात}\label{cmp:thy:cmp:sec}

$A$ संगणनक्षमपणे प्रगणनीय आहे असे समजा. मग~$A$ चा पूरक संच,
$\Complement{A} = \Nat \setminus A$, हाही नेहमी संगणनक्षमपणे
प्रगणनीय असतो का? पुढील प्रमेय आणि अनुप्रमेय दाखवतात की उत्तर
``नाही'' असे आहे.

\begin{thm}
\label{cmp:thy:cmp:thm:ce-comp}
$A$ हा नैसर्गिक संख्यांचा कोणताही संच असू द्या. $A$ संगणनक्षम असतो
तेव्हा आणि केवळ तेव्हाच, जेव्हा $A$ आणि $\Complement{A}$ हे दोन्ही
संगणनक्षमपणे प्रगणनीय असतात.
\end{thm}

\begin{proof}
एका दिशेची सिद्धता सोपी आहे: $A$ संगणनक्षम असेल, तर
$\Complement{A}$ देखील संगणनक्षम आहे ($\Char{A} = 1 \tsub
\Char{\Complement{A}}$). म्हणून दोन्ही संच संगणनक्षमपणे प्रगणनीय आहेत.

उलट दिशेने, $A$ आणि~$\Complement{A}$ हे दोन्ही संगणनक्षमपणे
प्रगणनीय आहेत असे समजा. $A$ हे~$\cfind{d}$ चे परिभाषाक्षेत्र
आणि $\Complement{A}$ हे~$\cfind{e}$ चे परिभाषाक्षेत्र असू द्या.
$h$ ची व्याख्या अशी करा:
\[h(x) = \umin{s}{(T(d,x,s) \lor T(e,x,s))}.
\]
म्हणजे, आदान~$x$ वर $h$ हे~$\cfind{d}$ किंवा~$\cfind{e}$
यांचे थांबणारे संगणन शोधते. $x \in A$ असेल, तर पहिल्या
पर्यायात शोध यशस्वी होतो; $x \in
\Complement{A}$ असेल, तर दुसऱ्या पर्यायात. म्हणून $h$
हे सर्वत्र परिभाषित संगणनक्षम फलन आहे. आता प्रत्येक~$x$
साठी $x \in A$ तेव्हा आणि केवळ तेव्हाच, जेव्हा $T(d, x, h(x))$, म्हणजे
या आदानावर $\cfind{d}$ परिभाषित असते. $T(d, x, h(x))$
हा संगणनक्षम संबंध असल्यामुळे $A$ संगणनक्षम आहे.

\end{proof}

\begin{explain}
अनौपचारिक संगणकीय भाषेत ही गोष्ट अधिक सोपी आहे: $A$
निर्णीत करण्यासाठी आदान $x$ वर $\cfind{d}$ आणि $\cfind{e}$
यांच्या थांबणाऱ्या संगणनांचा शोध घ्या. त्यांपैकी एक नक्की थांबते;
$\cfind{d}$ थांबले, तर $x$ हा~$A$ मध्ये आहे; अन्यथा $x$
हा~$\Complement{A}$ मध्ये आहे.
\end{explain}

\begin{cor}
\label{cmp:thy:cmp:cor:comp-k}
$\Complement{K_0}$ संगणनक्षमपणे प्रगणनीय नाही.
\end{cor}

\begin{proof}
$K_0$ संगणनक्षमपणे प्रगणनीय आहे, पण संगणनक्षम नाही, हे
आपल्याला माहीत आहे. $\Complement{K_0}$ संगणनक्षमपणे प्रगणनीय
असता, तर \ref{cmp:thy:cmp:thm:ce-comp} मुळे $K_0$ संगणनक्षम ठरला असता;
हे \ref{cmp:thy:ncp:thm:K-0} ला विरोधी आहे.
\end{proof}











\section{न्यूनीकरणक्षमता}\label{cmp:thy:red:sec}

\begin{explain}
आता आपल्याला माहीत आहे की $K_0$ हा किमान एक संच संगणनक्षमपणे
प्रगणनीय आहे, पण संगणनक्षम नाही. असे इतरही संच असतील हे उघड
आहे. प्रत्येक वेळी मूलभूत व्याख्यांपासून पुन्हा सुरुवात न करता
इतर संचांचे असे गुणधर्म सिद्ध करण्याची प्रभावी पद्धत
न्यूनीकरणक्षमता देते.

सर्वसाधारणपणे, $A$ या संचाचे~$B$ कडे ``न्यूनीकरण'' म्हणजे
घटक हे~$B$ मध्ये आहेत की नाहीत या प्रश्नांची उत्तरे
रूपांतरित करून घटक हे~$A$ मध्ये आहेत की नाहीत या प्रश्नांची उत्तरे
मिळवण्याची पद्धत. येथे आपण ``अनेक-एक न्यूनीकरणक्षमता''
विचारात घेऊ; वेगवेगळे गुणधर्म असणाऱ्या न्यूनीकरणक्षमतेच्या इतरही
अनेक संकल्पना आहेत. संगणकीय गुंतागुंतीच्या अभ्यासातही अशा
संकल्पना महत्त्वाच्या असतात; तेथे कार्यक्षमतेचाही विचार करावा
लागतो. उदाहरणार्थ, एखादा संच NP मध्ये असेल आणि प्रत्येक
NP समस्या त्याच्याकडे न्यूनीत करता येत असेल, तर त्याला
``NP-पूर्ण'' म्हणतात. येथे पुढे येणाऱ्या न्यूनीकरणासारखेच
न्यूनीकरण तेथे वापरतात, पण त्यासाठी ते बहुपदी वेळेत संगणित
करता येण्याची अतिरिक्त अट असते.

न्यूनीकरणाची कल्पना आपण आधीच अप्रत्यक्षपणे वापरली आहे.
$K$ हा संच पुढीलप्रमाणे ठरवा:
\[K = \Setabs{x}{\cfind{x}(x) \fdefined},
\]
म्हणजे $K = \Setabs{x}{x \in W_x}$. थांबण्याची समस्या
अनिर्णेय असल्याचा आपला पुरावा (\ref{cmp:thy:hlt:thm:halting-problem})
सर्वांत थेटपणे $K$ संगणनक्षम नाही हे दाखवतो. $K_0$ हा
पुढील संच आहे, हे आठवा:
\[K_0 = \Setabs{\tuple{e, x}}{\cfind{e}(x) \fdefined },
\]
म्हणजे $K_0 = \Setabs{\tuple{e,x}}{x \in W_e}$.

$K$ च्या असंगणनक्षमतेच्या कोणत्याही पुराव्याचा विस्तार करून
$K_0$ ची असंगणनक्षमता दाखवणे सोपे आहे: $K_0$ संगणनक्षम असता,
तर घटक~$x$ हा $K$ मध्ये आहे की नाही हे फक्त
$\tuple{x, x}$ ही जोडी $K_0$ मध्ये आहे का ते विचारून
निर्णीत करता आले असते. $x$ ला $\tuple{x, x}$ कडे नेणारे
फलन~$f$ हे $K$ चे $K_0$ कडे \emph{न्यूनीकरण} करण्याचे
उदाहरण आहे.
\end{explain}

\begin{defn}
$A$ आणि $B$ हे नैसर्गिक संख्यांचे संच असू द्या.
संगणनक्षम फलन~$f\colon \Nat \to \Nat$ हे $A$ चे~$B$
कडे \emph{अनेक-एक न्यूनीकरण} असते जर आणि फक्त जर
प्रत्येक नैसर्गिक संख्या~$x$ साठी
\[x \in A \quad \text{तेव्हा आणि केवळ तेव्हाच} \quad f(x) \in B.
\]
असे न्यूनीकरण~$f$ अस्तित्वात असल्यास, $A$ हे~$B$ कडे
\emph{अनेक-एक न्यूनीत करता येते} असे म्हणतो; हे
$A \leq_m B$ असे लिहितो. $A$ चे~$B$ कडे आणि उलटही
अनेक-एक न्यूनीकरण शक्य असल्यास, $A$ आणि $B$ हे
\emph{अनेक-एक सममूल्य} आहेत असे म्हणतो; हे
$A \equiv_m B$ असे लिहितो.
\end{defn}

वरील व्याख्येतील फलन $f$ एकास-एक असेल, तर $A$ हे~$B$
कडे \emph{एकास-एक न्यूनीत करता येते} असे म्हणतात.
पुढे दिलेली बहुतेक न्यूनीकरणे ही अधिक मजबूत अट पूर्ण करतात;
पण आपण ती बाब वापरणार नाही.

\begin{digress}
एकास-एक न्यूनीकरणक्षमता ही अनेक-एक न्यूनीकरणक्षमतेपेक्षा खरोखर
अधिक मजबूत अट आहे; हे खरे आहे, पण अजिबात उघड नाही. म्हणजे,
असे अनंत संच $A$ आणि~$B$ आहेत की $A$ हे~$B$ कडे अनेक-एक
न्यूनीत करता येते, पण~$B$ कडे एकास-एक न्यूनीत करता येत नाही.
\end{digress}











\section{न्यूनीकरणक्षमतेचे गुणधर्म}\label{cmp:thy:ppr:sec}

$A$ चे~$B$ कडे न्यूनीकरण होत असेल, तर आपण $A \leq_m B$
असे लिहितो. या चिन्हांकनावरून असे सुचते की $A \leq_m B$
असल्यास $A$ हे~$B$ पेक्षा ``अधिक कठीण नाही'' आणि $B$
हे~$A$ इतके किंवा त्याहून ``अधिक कठीण'' आहे. तसेच नेहमीचे
संख्यांवरील $\le$ जसे क्रम लावते, तसे $\le_m$ संचांना
(किंवा निर्णय समस्यांना) त्यांच्या गुंतागुंतीनुसार क्रम लावते.
पुढील दोन विधाने ही अंतःप्रेरणा समर्थित करतात. पहिले विधान
$\le_m$ संक्रमक आहे असे सांगते.

\begin{prop}
\label{cmp:thy:ppr:prop:trans-red}
$A \leq_m B$ आणि $B \leq_m C$ असल्यास $A \leq_m C$.
\end{prop}

\begin{proof}
$A$ चे~$B$ कडे न्यूनीकरण आणि $B$ चे~$C$ कडे न्यूनीकरण
यांचे संयोजन केल्यास $A$ चे~$C$ कडे न्यूनीकरण मिळते.
\end{proof}

\begin{prob}
$f$ आणि~$g$ ही अनुक्रमे $A$ चे~$B$ कडे आणि $B$ चे~$C$
कडे अनेक-एक न्यूनीकरणे असतील, तर $\comp{f}{g}$ हे $A$
चे~$C$ कडे अनेक-एक न्यूनीकरण आहे, हे दाखवून
\ref{cmp:thy:ppr:prop:trans-red} सिद्ध करा.
\end{prob}

\begin{prop}
\label{cmp:thy:ppr:prop:reduce}
$A$ आणि $B$ हे कोणतेही संच असू द्या आणि $A \leq_m B$
असे समजा.
\begin{enumerate}
\item $B$ संगणनक्षमपणे प्रगणनीय असेल, तर $A$ देखील तसाच आहे.
\item $B$ संगणनक्षम असेल, तर $A$ देखील तसाच आहे.
\end{enumerate}
\end{prop}

\begin{proof}
$f$ हे $A$ चे~$B$ कडे अनेक-एक न्यूनीकरण असू द्या.
पहिल्या विधानासाठी, $B$ हे आंशिक फलन~$g$ चे परिभाषाक्षेत्र
असेल, तर $A$ हे~$\comp{f}{g}$ चे परिभाषाक्षेत्र आहे हे
तपासा:
\begin{align*}
x \in A & \text{ तेव्हा आणि केवळ तेव्हाच } f(x) \in B \\
& \text{ तेव्हा आणि केवळ तेव्हाच }  g(f(x)) \fdefined.
\end{align*}

दुसऱ्या विधानासाठी, $B$ संगणनक्षम असेल, तर $B$
आणि~$\Complement{B}$ हे दोन्ही संगणनक्षमपणे प्रगणनीय
असतात (\ref{cmp:thy:cmp:thm:ce-comp}), हे आठवा. $f$ हे
$\Complement{A}$ चे $\Complement{B}$ कडेही अनेक-एक
न्यूनीकरण आहे, हे तपासणे कठीण नाही. त्यामुळे या सिद्धतेच्या
पहिल्या भागानुसार $A$ आणि~$\Complement{A}$ हे
संगणनक्षमपणे प्रगणनीय आहेत. म्हणून \ref{cmp:thy:cmp:thm:ce-comp} मुळे
$A$ देखील संगणनक्षम आहे. (पर्यायाने,
$\Char{A} = \comp{f}{\Char{B}}$ हे तपासा; मग
$\Char{B}$ संगणनक्षम असेल, तर~$\Char{A}$ देखील आहे.)
\end{proof}

\begin{prob}
$f$ हे $A$ चे~$B$ कडे अनेक-एक न्यूनीकरण आहे असे
समजा. मग $f$ हे $\Complement{A}$ चे $\Complement{B}$
कडेही अनेक-एक न्यूनीकरण आहे हे दाखवा.
\end{prob}

\begin{prob}
$f\colon \Nat \to \Nat$ हे $A$ चे~$B$ कडे अनेक-एक न्यूनीकरण असेल, तर
$\Char{A}
= \comp{f}{\Char{B}}$ हे दाखवा.

\end{prob}

\begin{digress}
\emph{ट्यूरिंग न्यूनीकरणक्षमता} ही न्यूनीकरणक्षमतेची अधिक
सर्वसाधारण संकल्पना इतर संदर्भांत, विशेषतः अनिर्णेयतेचे निष्कर्ष
सिद्ध करण्यासाठी, उपयोगी आहे. \ref{cmp:thy:cmp:cor:comp-k} नुसार
$K_0$ चा पूरक संच~$K_0$ कडे न्यूनीत करता येत नाही, कारण
तो संगणनक्षमपणे प्रगणनीय नाही. पण अंतःप्रेरणेने, $K_0$
विषयीच्या प्रश्नांची उत्तरे माहीत असतील, तर त्याच्या पूरकाविषयीची
उत्तरेही माहीत असतील. संगणनक्षम प्रक्रिया~$B$ विषयी प्रश्न
विचारून~$A$ विषयीच्या प्रश्नांची उत्तरे ठरवू शकत असेल, तर
$A$ हे~$B$ कडे ट्यूरिंग न्यूनीत करता येते असे म्हणतात.
अनेक-एक न्यूनीकरणक्षमतेपेक्षा ही अधिक व्यापक संकल्पना आहे:
अनेक-एक न्यूनीकरणात (1)~$B$ विषयी एकच प्रश्न विचारता येतो,
आणि (2)~``होय'' या उत्तराचे $A$ विषयीच्या प्रश्नासाठी
``होय'' असेच, तसेच ``नाही'' चे ``नाही'' असेच रूपांतर व्हावे
लागते. $A$ हे~$B$ कडे ट्यूरिंग न्यूनीत करता येत असेल
आणि $B$ संगणनक्षम असेल, तर $A$ देखील संगणनक्षम असतो;
परंतु आपण पाहिल्याप्रमाणे, संगणनक्षमपणे प्रगणनीय असण्याबाबत
असेच विधान खरे ठरत नाही.

आपण चर्चिलेल्या न्यूनीकरणक्षमतेच्या वेगवेगळ्या संकल्पनांचा
आणि त्यांतील फरकांचा विचार करा. मात्र या प्रकरणात आपण फक्त
अनेक-एक न्यूनीकरणक्षमता वापरू. आधीच्या परिच्छेदातील दोन्ही
प्रकारांना संगणकीय गुंतागुंतीत समांतर संकल्पना आहेत; त्यांत
ट्यूरिंग यंत्रे बहुपदी वेळेत चालण्याची अतिरिक्त अट असते.
अनेक-एक न्यूनीकरणक्षमतेचे गुंतागुंत-संबंधी रूप
\emph{कार्प न्यूनीकरणक्षमता}, तर ट्यूरिंग न्यूनीकरणक्षमतेचे
तसे रूप \emph{कुक न्यूनीकरणक्षमता} म्हणून ओळखले जाते.
\end{digress}











\section{संपूर्ण संगणनक्षमपणे प्रगणनीय संच}\label{cmp:thy:cce:sec}

\begin{defn}
संच $A$ हा (अनेक-एक न्यूनीकरणक्षमतेखाली) \emph{संपूर्ण
संगणनक्षमपणे प्रगणनीय संच} असतो, जर
\begin{enumerate}
\item $A$ संगणनक्षमपणे प्रगणनीय असेल, आणि
\item इतर कोणत्याही संगणनक्षमपणे प्रगणनीय संच $B$ साठी $B \leq_m A$ असेल.
\end{enumerate}
\end{defn}

म्हणजे, संपूर्ण संगणनक्षमपणे प्रगणनीय संच हे अशा संचांमधील
शक्य तितके ``कठीण'' संच असतात. त्यांच्याद्वारे
\emph{कोणत्याही} संगणनक्षमपणे प्रगणनीय संचाविषयीच्या
प्रश्नांची उत्तरे मिळवता येतात.

\begin{thm}
$K$, $K_0$, आणि $K_1$ हे तिन्ही संपूर्ण संगणनक्षमपणे
प्रगणनीय संच आहेत.
\end{thm}

\begin{proof}
$K_0$ संपूर्ण आहे हे पाहण्यासाठी $B$ हा कोणताही
संगणनक्षमपणे प्रगणनीय संच असू द्या. मग काही निर्देशांक $e$
साठी
\[B = W_e = \Setabs{x}{\cfind{e}(x) \fdefined}.
\]
$f(x) = \tuple{e, x}$ हे फलन $f$ असू द्या. मग प्रत्येक
नैसर्गिक संख्या $x$ साठी $x \in B$ तेव्हा आणि केवळ तेव्हाच, जेव्हा
$f(x) \in K_0$. म्हणजे $f$ हे $B$ चे~$K_0$ कडे
न्यूनीकरण करते.

$K_1$ संपूर्ण आहे हे पाहण्यासाठी, \ref{cmp:thy:k1:prop:k1}
च्या सिद्धतेत आपण $K_0$ चे त्याच्याकडे न्यूनीकरण केले
आहे, हे लक्षात घ्या. म्हणून \ref{cmp:thy:ppr:prop:trans-red}
नुसार कोणताही संगणनक्षमपणे प्रगणनीय संच~$K_1$ कडेही
न्यूनीत करता येतो.

$K_0$ चे~$K$ कडेही अशाच प्रकारे न्यूनीकरण करता येते.

\end{proof}

\begin{prob}
$K$ चे $K_0$ कडे न्यूनीकरण द्या.
\end{prob}

\begin{digress}
आतापर्यंत आपण पाहिलेल्या संगणनक्षमपणे प्रगणनीय संचांची सर्व
उदाहरणे एकतर संगणनक्षम किंवा संपूर्ण आहेत. हे विचित्र वाटले
पाहिजे!{} संगणनक्षमही नाहीत आणि संपूर्णही नाहीत, अशी
संगणनक्षमपणे प्रगणनीय संचांची उदाहरणे आहेत का? उत्तर होय
आहे; मात्र 1950 च्या दशकाच्या मध्यापर्यंत हे सिद्ध झाले
नव्हते. फ्रीडबर्ग आणि मुछ्निक यांनी स्वतंत्रपणे ते सिद्ध केले.
\end{digress}











\section{न्यूनीकरणक्षमतेचे एक उदाहरण}\label{cmp:thy:k1:sec}

\ref{cmp:thy:ppr:prop:reduce} चा एक उपयोग पाहू.

\begin{prop}
\label{cmp:thy:k1:prop:k1}
पुढीलप्रमाणे घ्या:
\[K_1 = \Setabs{e}{\cfind{e}(0) \fdefined}.
\]
मग $K_1$ संगणनक्षमपणे प्रगणनीय आहे, पण संगणनक्षम नाही.
\end{prop}

\begin{proof}
$K_1 = \Setabs{e}{\lexists[s][T(e,0,s)]}$ असल्याने,
\ref{cmp:thy:eqc:thm:exists-char} नुसार $K_1$ संगणनक्षमपणे
प्रगणनीय आहे.

$K_1$ संगणनक्षम नाही हे दाखवण्यासाठी $K_0$ चे
त्याच्याकडे न्यूनीकरण होते हे दाखवू.

\begin{explain}
हे थोडे गुंतागुंतीचे आहे, कारण $K_1$ वापरून आपण फक्त
विशिष्ट आदान $0$ पासून सुरू होणाऱ्या संगणनांबद्दल प्रश्न
विचारू शकतो. अशा प्रश्नांची उत्तरे देणारा हुशार मित्र
तुमच्याकडे आहे असे समजा (अशा मित्रांना ``ओरॅकल'' म्हणतात).
आता कोणीतरी तुम्हाला $\tuple{e,
  x}$ ही जोडी $K_0$ मध्ये आहे का, म्हणजे यंत्र $e$ हे
आदान $x$ वर थांबते का, असे विचारते. तुम्ही दुसरे यंत्र,
$e_x$, तयार करू शकता. ते \emph{कोणतेही} आदान मिळाले तरी
त्याकडे दुर्लक्ष करते आणि त्याऐवजी~$e$ हे यंत्र आदान
$x$ वर चालवते. मग यंत्र $e$ आदान $x$ वर थांबते का हा
प्रश्न, यंत्र $e_x$ आदान $0$ वर (किंवा इतर कोणत्याही
आदानावर) थांबते का या प्रश्नाशी समतुल्य आहे. आता नवीन
यंत्र $e_x$ आदान $0$ वर थांबते का हे मित्राला विचारा;
बदललेल्या प्रश्नाचे त्याचे उत्तर मूळ प्रश्नाचे उत्तरही देते.
यातून $K_0$ चे~$K_1$ कडे हवे ते न्यूनीकरण मिळते.
\end{explain}

सार्वत्रिक आंशिक संगणनक्षम फलन वापरून $f$ हे पुढीलप्रमाणे
व्याख्यायित त्रिस्थानी फलन असू द्या:
\[f(x,y,z) \simeq \cfind{x}(y).
\]
$f$ आपल्या तिसऱ्या आदानाकडे पूर्णपणे दुर्लक्ष करते, हे
लक्षात घ्या. $f = \cfind{e}[3]$ होईल असा निर्देशांक $e$
निवडा; म्हणजे
\[\cfind{e}[3](x,y,z) \simeq \cfind{x}(y).
\]
$s$-$m$-$n$ प्रमेयानुसार असे फलन $s(e,x,y)$ अस्तित्वात
आहे की प्रत्येक $z$ साठी
\begin{align*}
\cfind{s(e,x,y)}(z) & \simeq \cfind{e}[3](x,y,z) \\
& \simeq \cfind{x}(y).
\end{align*}

\begin{explain}
वरील अनौपचारिक युक्तिवादाच्या दृष्टीने $s(e,x,y)$ हा
अशा यंत्राचा निर्देशांक आहे की ते कोणतेही आदान $z$
मिळाले तरी त्याकडे दुर्लक्ष करते आणि $\cfind{x}(y)$
संगणित करते.
\end{explain}

विशेषतः, आपल्याला
\[\cfind{s(e,x,y)}(0) \fdefined \quad \text{तेव्हा आणि केवळ तेव्हाच} \quad
\cfind{x}(y) \fdefined.
\]
म्हणजे $\tuple{x, y} \in K_0$ तेव्हा आणि केवळ तेव्हाच, जेव्हा $s(e,x,y) \in
K_1$. कुठेही परिभाषित नसलेल्या आंशिक संगणनक्षम फलनाचा निर्देशांक $u$
निश्चित करा. मग पुढीलप्रमाणे व्याख्यायित फलन $g$:
\[g(w) = \begin{cases}
s(e,(w)_0,(w)_1) & \text{जर $w = \tuple{(w)_0,(w)_1}$ असेल तर} \\
u & \text{अन्यथा.}
\end{cases}
\]
हे $K_0$ चे~$K_1$ कडे न्यूनीकरण आहे. जोडीचा वैध संकेतांक असल्याची कसोटी
आदिम पुनरावर्ती आहे. अवैध संकेतांक $K_0$ मध्ये नसतो आणि त्याला $K_1$
मध्ये नसलेला निर्देशांक $u$ मिळतो; वैध जोड्यांसाठी वरील सममूल्यत्व लागू होते.

\end{proof}











\section{सर्वत्र परिभाषितता अनिर्णेय आहे}\label{cmp:thy:tot:sec}

एखादी गोष्ट असंगणनक्षम आहे हे दाखवण्यासाठी $s$-$m$-$n$
प्रमेय वापरण्याचे आणखी एक उदाहरण पाहू. $\fn{Tot}$ हा
सर्वत्र परिभाषित संगणनक्षम फलनांच्या निर्देशांकांचा संच असू
द्या, म्हणजे
\[\fn{Tot} = \Setabs{x}{\text{प्रत्येक $y$ साठी, $\cfind{x}(y)\fdefined$}}.
\]

\begin{prop}
\label{cmp:thy:tot:prop:total}
$\fn{Tot}$ संगणनक्षम नाही.
\end{prop}

\begin{proof}
$\fn{Tot}$ संगणनक्षम नाही हे पाहण्यासाठी $K$ चे त्याच्याकडे
न्यूनीकरण होते हे दाखवणे पुरेसे आहे. $h(x,y)$ ची व्याख्या
पुढीलप्रमाणे करा:
\[h(x,y) \simeq
\begin{cases}
0 & \text{जर $x \in K$ असेल तर} \\
\fundefined & \text{अन्यथा}
\end{cases}
\]
$h(x,y)$ हे $y$ वर अजिबात अवलंबून नाही, हे लक्षात घ्या.
$h$ हे आंशिक संगणनक्षम आहे हे पाहणे कठीण नसावे: आदान
$x, y$ वर~$h$ संगणित करताना आधी~$\cfind{x}$ चे
आदान~$x$ वरील अनुकरण करतो; हे संगणन थांबले, तर
$h(x,y)$ हे $0$ परत करून थांबते. म्हणून $h(x,y)$ हे
$\Zero(\umin{s}{T(x,x,s)})$ इतकेच आहे; येथे $\Zero$
हे स्थिर शून्य फलन आहे.

$s$-$m$-$n$ प्रमेयानुसार असे आदिम पुनरावर्ती फलन~$k(x)$
अस्तित्वात आहे की प्रत्येक $x$ आणि~$y$ साठी
\[\cfind{k(x)}(y) =
\begin{cases}
0 & \text{जर $x \in K$ असेल तर} \\
\fundefined & \text{अन्यथा}
\end{cases}
\]
म्हणून $x \in K$ असल्यास $\cfind{k(x)}$ सर्वत्र परिभाषित
असते; अन्यथा ते अपरिभाषित असते. अशा प्रकारे $k$ हे $K$
चे~$\fn{Tot}$ कडे न्यूनीकरण आहे.
\end{proof}

\begin{digress}
खरे तर $\fn{Tot}$ संगणनक्षमपणे प्रगणनीयही नाही; त्याची
गुंतागुंत ``अंकगणितीय पदानुक्रमात'' अधिक वरच्या स्तरावर
येते. मात्र येथे आपण या अधिक मजबूत निष्कर्षाचा विचार करणार नाही.
\end{digress}











\section{राइसचे प्रमेय}\label{cmp:thy:rce:sec}

विचार केल्यास लक्षात येईल की \ref{cmp:thy:tot:prop:total}
च्या सिद्धतेत $\fn{Tot}$ चे विशिष्ट गुणधर्म वापरलेले नाहीत.
$h(x,y)$ हे स्थिर फलन $j(y) = 0$ प्रमाणे वागते तेव्हा आणि केवळ
तेव्हाच, जेव्हा $x$ हा $K$ मध्ये असतो, अशी रचना आपण केली; पण त्या
परिस्थितीत ते इतर कोणत्याही आंशिक संगणनक्षम फलनाप्रमाणे
वागेल अशीही रचना करता आली असती. या निरीक्षणामुळे अधिक
सर्वसाधारण प्रमेय मांडता येते. साधारणपणे ते सांगते की
संगणनक्षम फलनांचा कोणताही अतुच्छ गुणधर्म निर्णेय नसतो.

$\cfind{0}$, $\cfind{1}$, $\cfind{2}$,~\dots हे आंशिक
संगणनक्षम फलनांचे आपले प्रमाण प्रगणन आहे, हे लक्षात ठेवा.

\begin{thm}[राइसचे प्रमेय]
  $C$ हा आंशिक संगणनक्षम फलनांचा कोणताही संच असू द्या आणि $A =
  \Setabs{n}{\cfind{n} \in C}$ असू द्या. $A$ संगणनक्षम असेल,
  तर $C$ रिकामा असतो किंवा $C$ हा सर्व आंशिक संगणनक्षम
  फलनांचा संच असतो.
\end{thm}

{\em निर्देशांकसंच} म्हणजे असा संच $A$ की $n$ आणि $m$ हे
एकच फलन ``संगणित'' करणारे निर्देशांक असतील, तर $n$ आणि
$m$ हे दोन्ही $A$ मध्ये असतात किंवा दोन्ही नसतात. प्रमेयातील संच~$A$
हा गुणधर्म पूर्ण करतो हे पाहणे कठीण नाही. उलट, $A$
निर्देशांकसंच असेल आणि $C$ हा त्या निर्देशांकांनी संगणित
होणाऱ्या फलनांचा संच असेल, तर $A = \Setabs{n}{\cfind{n} \in C}$.

\begin{explain}
या संज्ञेत राइसचे प्रमेय म्हणजे कोणताही अतुच्छ निर्देशांकसंच
निर्णेय नसतो असे विधान. प्रमेयाचा अर्थ समजण्यासाठी
\emph{कार्यक्रम} (उदा., तुमच्या आवडत्या प्रोग्रामिंग भाषेतील)
आणि ते संगणित करत असलेली फलने यांतील फरक लक्षात घ्या.
कार्यक्रम हे विन्यासात्मक वस्तू आहेत, आणि त्यांच्याविषयीचे
काही प्रश्न नक्कीच संगणनक्षम आहेत: या कार्यक्रमात 150
पेक्षा अधिक चिन्हे आहेत का? त्यात 22 पेक्षा अधिक ओळी
आहेत का? त्यात ``while'' विधान आहे का? ``print'' विधानाचे
आदान म्हणून ``hello world'' ही चिन्हमाला कुठे येते का?
राइसचे प्रमेय सांगते की कार्यक्रमाच्या \emph{वर्तनाविषयीचा}
कोणताही अतुच्छ प्रश्न संगणनक्षम नसतो. उदाहरणार्थ: हा
कार्यक्रम आदान $0$ वर थांबतो का? तो कधी तरी थांबतो का?
तो कधी सम संख्या परत करतो का?
\end{explain}

\begin{proof}[राइसच्या प्रमेयाची सिद्धता]
$C$ रिकामा नाही आणि सर्व आंशिक संगणनक्षम फलनांचा संचही
नाही असे समजा. $A$ हा~$C$ मधील फलनांच्या निर्देशांकांचा
संच असू द्या. $A$ संगणनक्षम असता, तर थांबण्याची समस्या
सोडवता आली असती, हे आपण दाखवू; म्हणून $A$ संगणनक्षम नाही.

व्यापकता न गमावता, कुठेही परिभाषित नसलेले फलन $f$ हे
$C$ मध्ये नाही असे धरू शकतो (तसे नसेल, तर पुढील
युक्तिवादात $C$ आणि त्याचा पूरक संच यांची अदलाबदल करा).
$g$ हे~$C$ मधील कोणतेही फलन असू द्या. कल्पना अशी आहे:
$A$ निर्णीत करता आले, तर~$f$ संगणित करणारे निर्देशांक
आणि~$g$ संगणित करणारे निर्देशांक यांतील फरक ओळखता
येईल; आणि ती क्षमता वापरून थांबण्याची समस्या सोडवता येईल.

कसे ते पाहू. सार्वत्रिक आंशिक संगणनक्षम फलने वापरून
पुढील फलन व्याख्यायित करता येते:
\[h(x,y) \simeq
\begin{cases}
\text{अपरिभाषित} & \text{जर $\cfind{x}(x) \fundefined$ असेल तर} \\
g(y) & \text{अन्यथा.}
\end{cases}
\]
$h$ संगणित करण्यासाठी आधी $\cfind{x}(x)$ संगणित करण्याचा
प्रयत्न करतो; ते संगणन थांबले, तर पुढे~$g(y)$ संगणित करतो;
आणि {\em हे} संगणन थांबले, तर त्याचे मूल्य परत करतो.
अधिक औपचारिकपणे,
\[h(x,y) \simeq \Proj{2}{0}(g(y),\fn{Un}(x,x)).
\]
असे लिहिता येते. येथे $\Proj{2}{0}(z_0, z_1) = z_0$ हे
$0$ वे आदान परत करणारे, संगणनक्षम $2$-स्थानी प्रक्षेपण
फलन आहे.

मग $h$ हे आंशिक संगणनक्षम फलनांचे संयोजन आहे. उजवी बाजू परिभाषित
असून तिचे मूल्य~$g(y)$ च्या समान असते तेव्हा आणि केवळ तेव्हाच, जेव्हा
$\fn{Un}(x,x)$ आणि $g(y)$ ही दोन्ही परिभाषित असतात.

निश्चित~$x$ साठी, $\cfind{x}(x)$ अपरिभाषित असेल, तर
प्रत्येक~$y$ साठी $h(x,y)$ अपरिभाषित असते; आणि
$\cfind{x}(x)$ परिभाषित असेल, तर $h(x,y) \simeq g(y)$.
म्हणून निश्चित~$x$ साठी $h(x,y)$ हे एकतर $f$ प्रमाणे
किंवा $g$ प्रमाणे वागते. $\cfind{x}(x)$ परिभाषित आहे
की नाही हे ठरवणे म्हणजे या दोन स्थितींपैकी कोणती आहे
हे ठरवणे. पण हेच $h_x(y) \simeq h(x,y)$ हे~$C$
मध्ये आहे की नाही हे ठरवण्यासारखे आहे; $A$ संगणनक्षम
असता, तर तसे करता आले असते.

अधिक औपचारिकपणे, $h$ आंशिक संगणनक्षम असल्याने तो काही
निर्देशांक~$e$ साठी $\cfind{e}$ या फलनाइतकाच आहे.
$s$-$m$-$n$ प्रमेयानुसार असे आदिम पुनरावर्ती फलन~$s$
आहे की प्रत्येक~$x$ साठी $\cfind{s(e,x)}(y) = h_x(y)$.
आता प्रत्येक $x$ साठी $\cfind{x}(x) \fdefined$ असेल,
तर $\cfind{s(e,x)}$ हे~$g$ इतकेच फलन आहे; म्हणून
$s(e,x)$ हा $A$ मध्ये आहे. उलट, $\cfind{x}(x)
\uparrow$ असेल, तर $\cfind{s(e,x)}$ हे~$f$ इतकेच फलन
आहे; म्हणून $s(e,x)$ हा~$A$ मध्ये नाही. दुसऱ्या शब्दांत,
प्रत्येक~$x$ साठी $x
\in K$ तेव्हा आणि केवळ तेव्हाच, जेव्हा $s(e,x) \in A$. $A$ संगणनक्षम
असता, तर~$K$ देखील संगणनक्षम ठरला असता; हा विरोधाभास
आहे. म्हणून $A$ संगणनक्षम नाही.
\end{proof}

राइसचे प्रमेय अतिशय प्रभावी आहे. पुढील तत्काळ अनुप्रमेयात
त्याच्या उपयोगांची काही उदाहरणे आहेत.
\begin{cor}
पुढील संच अनिर्णेय आहेत.
\begin{enumerate}
\item $\Setabs{x}{\text{$17$ हे $\cfind{x}$ च्या व्याप्तीत आहे}}$
\item $\Setabs{x}{\text{$\cfind{x}$ स्थिर आहे}}$
\item $\Setabs{x}{\text{$\cfind{x}$ सर्वत्र परिभाषित आहे}}$
\item $\Setabs{x}{\text{जेव्हा $y < y'$ असेल तेव्हा $\cfind{x}(y) \fdefined$, आणि
    जर $\cfind{x}(y') \fdefined$ असेल, तर $\cfind{x}(y) < \cfind{x}(y')$}}$
\end{enumerate}
\end{cor}

\begin{proof}
  हे सर्व अतुच्छ निर्देशांकसंच आहेत.
\end{proof}











\section{स्थिरबिंदू प्रमेय}\label{cmp:thy:fix:sec}

थांबण्याच्या समस्येचा पुन्हा विचार करू. तात्पुरत्या चिन्हांकनासाठी,
$\tuple{x, y}$ ऐवजी $\gn{\cfind{x}(y)}$ असे लिहू;
याला $\cfind{x}(y)$ या मूल्याचे ``नाव'' दर्शवणारे चिन्ह समजा.
या चिन्हांकनाने थांबण्याची समस्या अनिर्णेय असल्याचा आपला एक
पुरावा पुन्हा मांडता येतो.

प्रश्न: पुढील गुणधर्म असलेले संगणनक्षम फलन $h$ आहे का?
प्रत्येक $x$ आणि $y$ साठी,
\[h(\gn{\cfind{x}(y)}) =
\begin{cases}
1 & \text{जर $\cfind{x}(y) \fdefined$ असेल तर} \\
0 & \text{अन्यथा.}
\end{cases}
\]

उत्तर: नाही. अन्यथा पुढील आंशिक फलन
\[g(x) \simeq
\begin{cases}
0 & \text{जर $h(\gn{\cfind{x}(x)}) = 0$ असेल तर} \\
\text{अपरिभाषित} & \text{अन्यथा}
\end{cases}
\]
आंशिक संगणनक्षम असते आणि त्याला काही निर्देशांक~$e$ असता.
पण मग
\[\cfind{e}(e) \simeq
\begin{cases}
0 & \text{जर $h(\gn{\cfind{e}(e)}) = 0$ असेल तर} \\
\text{अपरिभाषित} & \text{अन्यथा,}
\end{cases}
\]
असे मिळते. त्यानुसार $\cfind{e}(e)$ परिभाषित असेल
तेव्हा आणि केवळ तेव्हाच ते अपरिभाषित असेल; हा विरोधाभास आहे.

आता $\cfind{e}$ असलेल्या समीकरणाकडे पाहा. त्यात एका
अर्थाने स्वसंदर्भ आहे: $\cfind{e}(e)$ चे मूल्य विशिष्ट
पद्धतीने $\gn{\cfind{e}(e)}$ वर अवलंबून राहील अशी
रचना केली आहे. स्थिरबिंदू प्रमेय सांगते की हे सर्वसाधारणपणे
करता {\em येते}—फक्त विरोधाभास सिद्ध करण्यासाठीच नाही.

\ref{cmp:thy:fix:lem:fixed-equiv} मध्ये स्थिरबिंदू प्रमेय मांडण्याच्या
दोन समतुल्य पद्धती दिल्या आहेत. तार्किकदृष्ट्या दोन्ही
विधाने खरी असल्याने ती समतुल्य आहेत असे म्हणता येते;
पण आपला खरा अर्थ असा आहे की प्रत्येक विधानावरून दुसरे
सहज निष्पन्न होते. म्हणून ती एकाच प्रमेयाची पर्यायी विधाने
मानता येतात.

\begin{lem}
\label{cmp:thy:fix:lem:fixed-equiv}
पुढील विधाने समतुल्य आहेत:
\begin{enumerate}
\item प्रत्येक आंशिक संगणनक्षम फलन $g(x,y)$ साठी असा
  निर्देशांक~$e$ आहे की प्रत्येक~$y$ साठी
\[\cfind{e}(y) \simeq g(e,y).
\]
\item प्रत्येक संगणनक्षम फलन~$f(x)$ साठी असा निर्देशांक~$e$
  आहे की प्रत्येक~$y$ साठी
\[\cfind{e}(y) \simeq \cfind{f(e)}(y).
\]
\end{enumerate}
\end{lem}

\begin{proof}
$(1) \Rightarrow (2)$: $f$ दिल्यावर $g(x,y) \simeq
\fn{Un}(f(x),y)$ असे $g$ व्याख्यायित करा. (1) वापरून
प्रत्येक~$y$ साठी पुढील गुणधर्म असलेला निर्देशांक~$e$
मिळवा:
\begin{align*}
\cfind{e}(y) & = \fn{Un}(f(e),y) \\
& = \cfind{f(e)}(y).
\end{align*}

$(2) \Rightarrow (1)$: $g$ दिल्यावर $s$-$m$-$n$
प्रमेय वापरून प्रत्येक $x$ आणि~$y$ साठी
$\cfind{f(x)}(y) \simeq g(x,y)$ होईल असे $f$ मिळवा.
(2) वापरून असा निर्देशांक~$e$ मिळवा की
\begin{align*}
\cfind{e}(y) & = \cfind{f(e)}(y) \\
& = g(e,y).
\end{align*}
याने सिद्धता पूर्ण झाली.
\end{proof}

\begin{explain}
विधान (1) खरे आहे (आणि म्हणून (2) देखील) हे दाखवण्यापूर्वी,
ते किती विलक्षण आहे याचा विचार करा. $e$ हा संगणक
कार्यक्रम आहे असे समजा; विधान (1) सांगते की कोणतेही
आंशिक संगणनक्षम $g(x,y)$ दिले असता, $g_e(y) \simeq g(e,y)$
संगणित करणारा संगणक कार्यक्रम $e$ शोधता येतो. म्हणजे
स्वतःचाच संदर्भ असलेले फलन संगणित करणारा कार्यक्रम शोधता येतो.
\end{explain}

\begin{thm}
\ref{cmp:thy:fix:lem:fixed-equiv} मधील दोन्ही विधाने खरी आहेत.
विशेषतः, प्रत्येक आंशिक संगणनक्षम फलन $g(x,y)$ साठी
असा निर्देशांक~$e$ आहे की प्रत्येक~$y$ साठी
\[\cfind{e}(y) \simeq g(e,y).
\]
\end{thm}

\begin{proof}
आवश्यक घटक वरच्या थांबण्याच्या समस्येवरील चर्चेत अप्रत्यक्षपणे
आहेतच. $\fn{diag}(x)$ हे असे संगणनक्षम फलन असू द्या की
ते प्रत्येक $x$ साठी $f_x(y) \simeq \cfind{x}(x,y)$
या फलनाचा निर्देशांक परत करते, म्हणजे
\[\cfind{\fn{diag}(x)}(y) \simeq \cfind{x}(x,y).
\]
$\fn{diag}$ कडे द्विस्थानी फलनाचा कार्यक्रम एकस्थानी
फलनाच्या कार्यक्रमात रूपांतरित करणारे फलन म्हणून पाहा;
मूळ कार्यक्रमाचा पहिला युक्तिवाद त्याच कार्यक्रमावर निश्चित
करून हे फलन मिळते. $\fn{diag}$ ची औपचारिक व्याख्या अशी
करता येते: आधी $s$ ची व्याख्या
\[s(x,y) \simeq \fn{Un}^2(x,x,y),
\]
अशी करा. येथे $\fn{Un}^2$ हे आंशिक
संगणनक्षम द्विस्थानी फलनांसाठी सार्वत्रिक असलेले त्रिस्थानी
फलन आहे. मग $s$-$m$-$n$ प्रमेयानुसार पुढील अट पूर्ण करणारे
आदिम पुनरावर्ती फलन $\fn{diag}$ मिळते:
\[\cfind{\fn{diag}(x)}(y) \simeq s(x,y).
\]

आता $l$ फलनाची व्याख्या अशी करा:
\[l(x,y) \simeq g(\fn{diag}(x),y).
\]
आणि $\gn{l}$ हा $l$ चा निर्देशांक असू द्या. शेवटी
$e = \fn{diag}(\gn{l})$ असू द्या. मग प्रत्येक $y$ साठी
\begin{align*}
\cfind{e}(y) & \simeq \cfind{\fn{diag}(\gn{l})}(y) \\
& \simeq \cfind{\gn{l}}(\gn{l}, y) \\
& \simeq l(\gn{l}, y) \\
& \simeq g(\fn{diag}(\gn{l}),y) \\
& \simeq g(e, y),
\end{align*}
हवे तसे.
\end{proof}

\begin{explain}
येथे नेमके काय घडते? स्वतःचाच मजकूर छापणारा संगणक
कार्यक्रम लिहायचा आहे असे समजा. यासाठी तुमची प्रोग्रामिंग
भाषा चिन्हमालांवरील समृद्ध व काहीशी विचित्र फलनांचा संच
पुरवते असेही समजा. विशेषतः, तिच्यात पुढीलप्रमाणे काम करणारे
फलन $\fn{diag}$ आहे असे समजा: आदान चिन्हमाला~$s$
दिली असता $\fn{diag}$ $s$ मधील `x' या चिन्हाच्या प्रत्येक
घटनेचा शोध घेते आणि त्या जागी मूळ चिन्हमालेची अवतरणचिन्हांत
घेतलेली आवृत्ती बसवते. उदाहरणार्थ, पुढील चिन्हमाला
\begin{quote}
\begin{verbatim}
hello x world
\end{verbatim}
\end{quote}
आदान म्हणून दिल्यास हे फलन पुढील चिन्हमाला परत करते:
\begin{quote}
\begin{verbatim}
hello 'hello x world' world
\end{verbatim}
\end{quote}
मग अपेक्षित कार्यक्रम लिहिणे सोपे आहे; पुढील कार्यक्रम
हे काम करतो हे तपासा:
\begin{quote}
\begin{verbatim}
print(diag('print(diag(x))'))
\end{verbatim}
\end{quote}
C++ आणि Java सारख्या अधिक प्रचलित प्रोग्रामिंग भाषांतही
हीच कल्पना (अधिक गुंतागुंतीच्या अंमलबजावणीसह) काम करते.

आता आपण स्थिरबिंदू प्रमेयाच्या सिद्धतेपासून काही पावलेच
दूर आहोत. $\fn{print}(x,y)$ हे print फलनाचे एक रूप
चिन्हमाला $x$ आणि संख्यात्मक युक्तिवाद $y$ घेते व
$x$ ही चिन्हमाला $y$ वेळा छापते असे समजा. मग पुढील ``कार्यक्रम''
\begin{quote}
\begin{verbatim}
getinput(y);
print(diag('getinput(y); print(diag(x), y)'), y)
\end{verbatim}
\end{quote}
आदान $y$ मिळाल्यावर स्वतःचा मजकूर $y$ वेळा छापतो.
$\fn{getinput}$---$\fn{print}$---$\fn{diag}$ या सांगाड्याच्या
जागी कोणतेही फलन $g(x,y)$ घेतल्यास
\begin{quote}
\begin{verbatim}
g(diag('g(diag(x), y)'), y)
\end{verbatim}
\end{quote}
मिळते. हा कार्यक्रम आदान $y$ वर स्वतःच्या कार्यक्रमाला
आणि $y$ ला $g$ चे युक्तिवाद म्हणून देतो. ``अवतरणचिन्हांत
घेणे'' म्हणजे ``निर्देशांक वापरणे'' असे मानल्यास वरील
सिद्धता मिळते.

सध्या हा पुरावा औपचारिक युक्ती किंवा जादू वाटला तरी
हरकत नाही. पण वरील युक्तिवादाचे तपशील पुन्हा उभारता
आले पाहिजेत. अपूर्णतेची प्रमेये (आणि संबंधित ``स्थिरबिंदू
प्रमेय'') सिद्ध करताना ते का काम करते हे समजण्याच्या
इतर पद्धतींवर चर्चा करू.
\end{explain}













\section{स्थिरबिंदू प्रमेयाचा उपयोग}\label{cmp:thy:apf:sec}

स्थिरबिंदू प्रमेयामुळे आंशिक संगणनक्षम फलनांची व्याख्या
त्यांच्या निर्देशांकांच्या साहाय्याने करता येते. उदाहरणार्थ,
प्रत्येक $y$ साठी पुढील अट पूर्ण करणारा निर्देशांक $e$
शोधता येतो:
\[\cfind{e}(y) = e + y.
\]
आणखी एक उदाहरण म्हणजे स्थिरबिंदू प्रमेयाची सिद्धता
वापरून Java किंवा C++ मध्ये स्वतःचा मजकूर छापणारा
कार्यक्रम तयार करता येतो.

प्रत्येक $e$ साठी $W_e$ हे $\cfind{e}$ चे परिभाषाक्षेत्र
असू द्या. मग $W_0$, $W_1$, $W_2$,~\dots ही क्रमिका
संगणनक्षमपणे प्रगणनीय संचांचे प्रगणन करते, हे आठवा.
यांतील काही संच संगणनक्षम आहेत. आदान म्हणून $x$
घेऊन, $W_x$ संगणनक्षम असेल तेव्हा त्याच्या जातिबोधक
फलनाचा निर्देशांक परत करणारा अल्गोरिदम आहे का, असा
प्रश्न विचारता येतो. उत्तर ``नाही'' आहे; असा अल्गोरिदम नाही:

\begin{thm}
पुढील गुणधर्म असलेले कोणतेही आंशिक संगणनक्षम फलन $f$
नाही: $W_e$ संगणनक्षम असेल तेव्हा $f(e)$ परिभाषित
असते आणि $\cfind{f(e)}$ हे त्याचे जातिबोधक फलन असते.
\end{thm}

\begin{proof}
$f$ हे कोणतेही आंशिक संगणनक्षम फलन असू द्या. आपण असा
निर्देशांक $e$ तयार करू की $W_e$ संगणनक्षम असेल, पण
अपेक्षित निर्देशांक मिळतच नसेल किंवा मिळालेले $\cfind{f(e)}$
हे त्याचे जातिबोधक फलन नसेल.

स्थिरबिंदू प्रमेय वापरून असा निर्देशांक $e$ मिळतो की
\[\cfind{e}(y) \simeq
\begin{cases}
0 & \text{जर $y=0$ आणि $\cfind{f(e)}(0) \fdefined = 0$ असेल तर} \\
\text{अपरिभाषित} & \text{अन्यथा.}
\end{cases}
\]
म्हणजे पुढीलप्रमाणे व्याख्यायित फलनावर स्थिरबिंदू प्रमेय
लावून $e$ मिळतो:
\[g(x,y) \simeq
\begin{cases}
0 & \text{जर $y=0$ आणि $\cfind{f(x)}(0) \fdefined = 0$ असेल तर} \\
\text{अपरिभाषित} & \text{अन्यथा.}
\end{cases}
\]
$g$ आंशिक संगणनक्षम आहे हे अनौपचारिकपणे असे पाहता
येते: आदान $x$ आणि $y$ मिळाल्यावर अल्गोरिदम आधी
$y$ हे~$0$ शी समान आहे का ते तपासतो. समान असल्यास
तो $f(x)$ संगणित करतो आणि मग सार्वत्रिक यंत्र वापरून
$\cfind{f(x)}(0)$ संगणित करतो. हे शेवटचे संगणन थांबून
$0$ परत करत असेल, तर अल्गोरिदम~$0$ परत करतो;
अन्यथा अल्गोरिदम थांबत नाही.

आता $\cfind{f(e)}(0)$ परिभाषित असून $0$ च्या समान
असेल, तर $\cfind{e}(y)$ परिभाषित असते तेव्हा आणि केवळ
तेव्हाच $y$~$0$ च्या समान असते; म्हणून $W_e =
\{ 0 \}$. $\cfind{f(e)}(0)$ अपरिभाषित असेल किंवा
परिभाषित असूनही $0$ च्या समान नसेल, तर $W_e = \emptyset$.
दोन्ही स्थितींत अपेक्षित गुणधर्म फसतो: या आदानावर फलन
निर्देशांक परत करतच नसेल, तर अपेक्षित परिभाषितता मिळत
नाही; निर्देशांक परत करत असेल, तर $\cfind{f(e)}$ हे~$W_e$
चे जातिबोधक फलन नाही, कारण आदान $0$ वर ते चुकीचे
उत्तर देते.
\end{proof}











\section{स्वसंदर्भ वापरून फलनांची व्याख्या}\label{cmp:thy:slf:sec}

फलनांची व्याख्या त्यांच्याच साहाय्याने करता येणे साधारणपणे
उपयुक्त असते. उदाहरणार्थ, संगणनक्षम फलने $k$, $l$
आणि~$m$ दिली असता, स्थिरबिंदू उपप्रमेयानुसार प्रत्येक~$y$
साठी पुढील समीकरण पूर्ण करणारे आंशिक संगणनक्षम फलन~$f$
अस्तित्वात असते:
\[f(y) \simeq
\begin{cases}
k(y) & \text{जर $l(y) = 0$ असेल तर} \\
f(m(y)) & \text{अन्यथा.}
\end{cases}
\]
अधिक नेमके सांगायचे, तर
\[g(x,y) \simeq
\begin{cases}
k(y) & \text{जर $l(y) = 0$ असेल तर} \\
\cfind{x}(m(y)) & \text{अन्यथा}
\end{cases}
\]
असे घेऊन, नंतर $\cfind{e}(y) = g(e,y)$ होईल असा
निर्देशांक~$e$ स्थिरबिंदू उपप्रमेयाने मिळवल्यावर~$f$ मिळते.

ठोस उदाहरण म्हणून, ``महत्तम साधारण विभाजक'' फलन
$\fn{gcd}(u,v)$ ची व्याख्या अशी करता येते:
\[\fn{gcd}(u,v) \simeq
\begin{cases}
v & \text{जर $u = 0$ असेल तर} \\
\fn{gcd}(\fn{mod}(v, u), u) & \text{अन्यथा}
\end{cases}
\]
येथे $\fn{mod}(v, u)$ म्हणजे $v$ ला~$u$ ने भागल्यावर
उरणारी बाकी. स्थिरबिंदू उपप्रमेय वापरून $\fn{gcd}$
आंशिक संगणनक्षम आहे हे मिळते. (प्रत्यक्षात, $y$ मध्ये
$\tuple{u, v}$ या जोडीचा संकेतांक घेऊन हे वरच्या
रूपात मांडता येते.) मग~$u$ वरील गणितीय विगमनाने $\fn{gcd}$
हे प्रत्यक्षात सर्वत्र परिभाषित आहे हे दिसते.

अर्थात, नुकत्याच पाहिलेल्या उदाहरणांपेक्षा बरीच अधिक
गुंतागुंतीची स्वसंदर्भात्मक व्याख्या करता येते. बहुतेक
प्रोग्रामिंग भाषा फलनांची व्याख्या त्यांच्याच साहाय्याने
करण्याचे कोणते ना कोणते साधन देतात. पण एखादे फलन
संगणित करणाऱ्या अल्गोरिदमच्या \emph{निर्देशांका}च्या
साहाय्याने त्याची व्याख्या करता येण्यापेक्षा हे काहीसे
कमी विलक्षण आहे. स्थिरबिंदू प्रमेय पूर्ण सर्वसाधारणतेने
नेमके तेच करू देते.



